\NormalFont\ShortTitle{Lucas}
{\MT Evangelho Segundo Lucas

\par }\ChapOne{1}{\PP \VerseOne{1}Como muitos empreenderam pôr em ordem o relato das coisas que entre nós se cumpriram,
\VS{2}conforme os que
{\ADD{as}} viram que desde o princípio, e foram servidores das palavra, nos transmitiram,
\VS{3}pareceu-me bom que também eu, que tenho me informado com exatidão desde o princípio, escrevesse-as em ordem a ti, caríssimo Teófilo,
\VS{4}para que conheças a certeza das coisas de que foste ensinado.
\VS{5}Houve nos dias de Herodes, rei da Judeia, um sacerdote chamado Zacarias, da ordem de Abias; sua mulher das filhas de Arão, e o seu nome era Isabel.
\VS{6}E ambos eram justos diante de Deus; andavam sem repreensão em todos os mandamentos e preceitos do Senhor.
\VS{7}Mas não tinham filho, porque Isabel era estéril, e ambos tinham idade avançada.
\FTNT{*}{{\FR 1:7 }
tinham idade avançada
lit. eram avançados em dias
}
\VS{8}E aconteceu que, quando ele fazia o trabalho sacerdotal diante de Deus, na ordem da sua turma,
\VS{9}conforme o costume do sacerdócio, foi sorteado a entrar no santuário do Senhor para oferecer o incenso.
\VS{10}E toda a multidão do povo estava fora orando, à hora do incenso.
\VS{11}Então um anjo do Senhor lhe apareceu em pé, à direita do altar do incenso.
\VS{12}Quando Zacarias
{\ADD{o}} viu, perturbou-se, e o medo o tomou.
\VS{13}Mas o anjo lhe disse: Zacarias, não temas, pois a tua oração foi ouvida. E Isabel, tua mulher, dará à luz um filho, e lhe porás o nome de João.
\VS{14}E terás prazer e alegria, e muitos se alegrarão do seu nascimento;
\VS{15}pois ele será grande diante do Senhor, e não beberá vinho nem bebida alcoólica, e será cheio do Espírito Santo, já desde o ventre da sua mãe.
\VS{16}E converterá muitos dos filhos de Israel ao Senhor seu Deus.
\VS{17}E irá adiante dele no espírito e virtude de Elias, para converter os corações dos pais aos filhos, e os rebeldes à prudência dos justos; a fim de preparar um povo pronto ao Senhor.
\VS{18}Então Zacarias disse ao anjo: Como terei certeza disso? Pois sou velho, e a minha mulher de idade avançada.
\FTNT{†}{{\FR 1:18 }
idade avançada
Lit. avançada em dias
}
\VS{19}O anjo lhe respondeu: Eu sou Gabriel, que fico presente diante de Deus, e fui enviado para te falar e te dar estas boas notícias.
\VS{20}Eis que ficarás mudo, e não poderás falar até o dia em que essas coisas aconteçam, porque não creste nas minhas palavras, que se cumprirão no devido tempo.
\VS{21}E o povo estava esperando Zacarias, e se surpreenderam de que ele demorava no santuário.
\VS{22}Quando ele saiu, não conseguia lhes falar; e entenderam que havia tido alguma visão no santuário. Ele lhes fazia gestos, e continuou mudo.
\VS{23}E sucedeu que, terminados os dias do seu serviço, voltou à sua casa.
\VS{24}E depois daqueles dias Isabel, sua mulher, Isabel, engravidou-se, e por cinco meses se escondeu, dizendo:
\VS{25}Pois isto o Senhor me fez nos dias em que ele
{\ADD{me}} observou, para acabar a minha humilhação entre as pessoas.
\VS{26}E no sexto mês o anjo Gabriel foi enviado por Deus a uma cidade da Galileia, chamada Nazaré,
\VS{27}a uma virgem prometida em casamento com um homem chamado José, da descendência
\FTNT{‡}{{\FR 1:27 }
Lit. “casa”
} de Davi; e o nome da virgem era Maria.
\VS{28}O anjo entrou onde ela estava, e disse: Alegra-te, agraciada; o Senhor
{\ADD{é}} contigo, bendita
{\ADD{és}} tu entre as mulheres.
\VS{29}Quando ela
{\ADD{o}} viu, perturbou-se muito por suas palavras, e se perguntava que saudação seria esta.
\VS{30}Então o anjo lhe disse: Maria, não temas, porque encontraste graça diante de Deus.
\VS{31}E eis que em teu ventre conceberás, darás à luz um filho, e chamarás o seu nome de Jesus.
\VS{32}Ele será grande, e será chamado Filho do Altíssimo; e o Senhor Deus lhe dará o trono de Davi, seu ancestral.
\FTNT{§}{{\FR 1:32 }
Lit. pai
}
\VS{33}E reinará eternamente na casa de Jacó; o seu reino não terá fim.
\VS{34}Maria disse ao anjo: Como será isso? Pois não conheço intimamente homem algum.
\VS{35}O anjo lhe respondeu: O Espírito Santo virá sobre ti, e o poder do Altíssimo te cobrirá com a sua sombra. Por isso o Santo que nascerá de ti será chamado Filho de Deus.
\VS{36}E eis que Isabel, tua prima, também está grávida de um filho na sua velhice; e este é o sexto mês para a que era chamada de estéril.
\VS{37}Pois para Deus nada será impossível.
\VS{38}Então Maria disse: Eis aqui a serva do Senhor; cumpra-se em mim conforme a tua palavra. E o anjo ausentou-se dela.
\VS{39}E naqueles dias, Maria levantou-se, e foi apressada à região montanhosa, a uma cidade de Judá.
\VS{40}Ela entrou na casa de Zacarias, e cumprimentou Isabel.
\VS{41}E aconteceu que, quando Isabel ouviu o cumprimento de Maria, a criança saltou no seu ventre, e Isabel foi cheia do Espírito Santo.
\VS{42}Então exclamou em alta voz: Bendita
{\ADD{és}} tu entre as mulheres, e bendito
{\ADD{é}} o fruto do teu ventre!
\VS{43}Como me acontece isto,
\FTNT{*}{{\FR 1:43 }
Lit. “De onde me acontece isto”
} que a mãe de meu Senhor venha a mim?
\VS{44}Pois eis que, quando a voz do teu cumprimento chegou aos meus ouvidos, a criança saltou de alegria no meu ventre.
\VS{45}E bendita
\FTNT{†}{{\FR 1:45 }
Ou: bem-aventurada
} é a que creu, pois as coisas que da parte do Senhor lhe foram ditas se cumprirão.
\VS{46}Então Maria disse: A minha alma engrandece ao Senhor,
\VS{47}e o meu espírito se alegra em Deus meu Salvador.
\VS{48}Porque ele observou a humilde condição da sua serva; pois eis que desde agora todas as gerações me chamarão bendita.
\FTNT{‡}{{\FR 1:48 }
Ou: bem-aventurada
}
\VS{49}Pois o Poderoso me fez grandes coisas; e santo
{\ADD{é}} o seu nome.
\VS{50}E a sua misericórdia é de geração em geração sobre os que o temem.
\VS{51}Com o seu braço manifestou poder; dispersou os soberbos no pensamento dos seus corações.
\VS{52}Derrubou os poderosos dos tronos, e elevou os humildes.
\VS{53}Encheu de bens aos famintos, e despediu vazios os ricos.
\VS{54}Auxiliou ao seu servo Israel, lembrando-se da
{\ADD{sua}} misericórdia,
\VS{55}como falou aos nossos ancestrais,
\FTNT{§}{{\FR 1:55 }
Lit. pais
} a Abraão, e à sua descendência,
\FTNT{*}{{\FR 1:55 }
Lit. semente
} para sempre.
\VS{56}Maria esteve com ela por quase três meses; então voltou para sua casa.
\VS{57}E completou-se à Isabel o tempo de dar à luz, e teve um filho.
\VS{58}Os seus vizinhos e parentes ouviram que o Senhor havia feito a ela grande misericórdia, e alegraram-se com ela.
\VS{59}E aconteceu que ao oitavo dia vieram circuncidar o menino, e o chamavam Zacarias, do nome do seu pai.
\VS{60}E sua mãe respondeu: Não, porém será chamado João.
\VS{61}E disseram-lhe: Ninguém há entre os teus parentes que se chame deste nome.
\VS{62}E perguntaram por gestos ao seu pai como queria que lhe chamassem.
\VS{63}Ele pediu uma tábua, e escreveu: O seu nome é João. E todos se surpeenderam.
\VS{64}E logo a sua boca se abriu, e a sua língua se
{\ADD{soltou}} ; e voltou a falar, louvando a Deus.
\VS{65}E veio temor sobre todos os seus vizinhos; e em toda a região montanhosa da Judeia todas essas coisas foram divulgadas.
\VS{66}E todos os que ouviam, guardavam em seus corações, dizendo: Quem será este menino? E a mão do Senhor estava com ele.
\VS{67}E Zacarias, seu pai, foi cheio do Espírito Santo, e profetizou, dizendo:
\VS{68}Bendito
{\ADD{seja}} o Senhor Deus de Israel, porque visitou e redimiu o seu povo.
\VS{69}E nos levantou uma poderosa salvação
\FTNT{†}{{\FR 1:69 }
Ou: “um poderoso Salvador”. Lit. “um chifre de salvação”
} na casa de seu servo Davi,
\VS{70}como falou pela boca dos seus santos profetas, desde o princípio do mundo;
\VS{71}Que
{\ADD{nos}} livraria dos nossos inimigos e da mão de todos os que nos odeiam,
\VS{72}para manifestar misericórdia aos nossos ancestrais,
\FTNT{‡}{{\FR 1:72 }
Lit. pais
} e se lembrar do seu santo pacto,
\VS{73}do juramento que fez a Abraão, nosso ancestral;
\VS{74}de nos conceder que, libertos da mão dos nossos inimigos, o serviríamos sem temor,
\VS{75}em santidade e justiça diante dele, todos os dias da nossa vida.
\VS{76}E tu, menino, serás chamado profeta do Altíssimo; porque irás adiante da face do Senhor, para preparar os seus caminhos,
\VS{77}para dar a seu povo conhecimento da salvação, na remissão de seus pecados;
\VS{78}por causa da profunda misericórdia
\FTNT{§}{{\FR 1:78 }
profunda misericórdia
Lit. das entranhas de misericórdia
} do nosso Deus, pela qual o raiar da manhã nos visitou;
\VS{79}para iluminar os que estão sentados nas trevas e sombra de morte; a fim de guiar os nossos pés pelo caminho da paz.
\VS{80}E o menino crescia e se fortalecia em espírito. E esteve nos desertos até o dia em que se mostrou a Israel.

\par }\Chap{2}{\PP \VerseOne{1}E aconteceu naqueles dias que saiu um decreto da parte de César Augusto, para que todo a terra se registrasse.
\VS{2}(Esse primeiro censo foi feito quando Quirino era o governador da Síria.)
\VS{3}E todos foram se registrar, cada um à sua própria cidade.
\VS{4}E José também subiu da Galileia, da cidade de Nazaré à Judeia, à cidade de Davi, que se chama Belém; (porque
{\ADD{ele}} era da casa e família de Davi.)
\VS{5}Para se registrar com Maria, sua mulher, com ele prometida em casamento, que estava grávida.
\VS{6}E aconteceu que, enquanto eles ali, completaram-se os dias em que ela havia de dar à luz.
\VS{7}E deu à luz seu filho primogênito, e o envolveu em panos, e o deitou numa manjedoura; porque não havia lugar para eles na hospedaria.
\VS{8}E naquela mesma localidade havia pastores que estavam no campo, e vigiavam o seu rebanho durante as vigílias da noite.
\VS{9}E eis que um anjo do Senhor lhe apareceu, e a glória do Senhor os cercou de resplendor. Então tiveram grande temor.
\VS{10}E o anjo lhes disse: Não temais; pois eis que vos dou notícias de grande alegria, que será para todo o povo:
\VS{11}que hoje na cidade de Davi vos nasceu o Salvador, que é Cristo, o Senhor.
\VS{12}E isto vos será por sinal: achareis o menino envolto em panos, deitado numa manjedoura.
\VS{13}E no mesmo instante apareceu com o anjo uma multidão de exércitos celestiais, louvando a Deus, e dizendo:
\VS{14}Glória nas alturas a Deus, paz na terra,
{\ADD{e}} aos seres humanos boa vontade.
\VS{15}E aconteceu que, quando os anjos se ausentaram deles para o céu, os pastores disseram uns aos outros: Vamos, pois, até Belém, e vejamos isso que aconteceu, e o Senhor nos informou.
\VS{16}Então foram apressadamente, e acharam Maria, e José, e o menino deitado na manjedoura.
\VS{17}Quando
{\ADD{o}} viram, contaram o que lhes havia sido dito sobre o menino.
\VS{18}E todos os que ouviram se admiraram com o que os pastores lhes diziam.
\VS{19}Mas Maria guardava todas essas palavras, considerando
{\ADD{-as}} no seu coração.
\VS{20}Os pastores voltaram glorificando e louvando a Deus por tudo o que haviam ouvido e visto, como lhes havia sido dito.
\VS{21}Quando se completaram os oito dias para circuncidar o menino, foi chamado o seu nome Jesus, o qual havia sido posto pelo anjo antes que fosse concebido.
\VS{22}E quando se completaram os dias da purificação dela, segundo a Lei de Moisés, o levaram a Jerusalém, para o apresentarem ao Senhor,
\VS{23}conforme o que está escrito na Lei do Senhor: Todo macho primogênito
\FTNT{*}{{\FR 2:23 }
primogênito
Lit. que abrir o ventre materno
} será consagrado ao Senhor.
{\RQ{Êxodo 13:2,12
}}
\VS{24}E para darem a oferta segundo o que se diz na Lei do Senhor: um par de rolinhas ou dois pombinhos.
\VS{25}E eis que havia em Jerusalém um homem de nome Simeão; e esse homem era justo e devoto, e esperava a consolação de Israel; e o Espírito Santo estava sobre ele.
\VS{26}E havia lhe sido revelado pelo Espírito Santo que ele não veria a morte antes de ter visto o Cristo do Senhor.
\VS{27}E pelo Espírito foi ao Templo; e quando os pais trouxeram o menino Jesus, para fazerem com ele conforme o costume da Lei,
\FTNT{†}{{\FR 2:27 }
a lei (ou instrução) de Moisés, a Torá
}
\VS{28}Então ele o tomou em seus braços, louvou a Deus, e disse:
\VS{29}Agora, Senhor, despedes o teu servo em paz, conforme a tua palavra,
\VS{30}pois os meus olhos viram a tua salvação,
\VS{31}a qual preparaste diante de todos os povos;
\VS{32}Luz para iluminar as gentios, e para a glória do teu povo Israel.
\VS{33}E José e sua mãe se admiraram das coisas que se diziam sobre ele.
\VS{34}Simeão os abençoou, e disse a Maria, sua mãe: Eis que este é posto para queda e levantamento de muitos em Israel; e para sinal que terá oposição,
\VS{35}(também uma espada atravessará a tua própria alma) para que os pensamentos de muitos corações se manifestem.
\VS{36}E estava ali a profetiza Ana, filha de Fanuel, da tribo de Aser. Ela já tinha idade avançada, e havia vivido com o marido sete anos desde sua virgindade.
\VS{37}E era viúva de cerca de oitenta e quatro anos, e não se afastava do Templo, servindo
{\ADD{a Deus}} com jejuns e orações, de noite e de dia.
\VS{38}E ela veio na mesma hora; ela agradeci ao Senhor, e falava dele a todos os que esperavam a redenção em Jerusalém.
\VS{39}E quando acabaram de cumprir tudo, segundo a Lei do Senhor, voltaram à Galileia, para a sua cidade de Nazaré.
\VS{40}E o menino crescia e se fortalecia em espírito, e cheio de sabedoria; e a graça de Deus estava sobre ele.
\VS{41}Todos os anos, seus pais iam a Jerusalém, para a festa da Páscoa.
\VS{42}E quando
{\ADD{Jesus}} tinha doze anos, subiram a Jerusalém, conforme o costume do festival;
\VS{43}E quando aqueles dias terminaram, eles partiram de volta, mas o menino Jesus ficou em Jerusalém sem que José e sua mãe soubessem.
\VS{44}Porém, como pensavam que ele vinha pelo caminho entre os companheiros de viagem, eles andaram o caminho de um dia; e o procuraram entre os parentes e conhecidos.
\VS{45}E como não o acharam, voltaram a Jerusalém em busca dele.
\VS{46}E aconteceu que, três dias depois, o acharam no templo, sentado no meio dos mestres, ouvindo-os, e perguntando-lhes.
\VS{47}E todos os que o ouviam admiravam a sua inteligência e respostas;
\VS{48}Quando eles o viram, ficarm surpresos. E sua mãe lhe disse: Filho, por que fizeste assim conosco? Eis que teu pai e eu com ansiedade te procurávamos.
\VS{49}Ele lhes disse: Por que me procuráveis? Não sabeis que devo estar nos negócios de meu Pai?
\VS{50}E eles não entenderam as palavras que lhes dizia.
\VS{51}Então desceu com eles, e veio a Nazaré, e os obedecia. E a sua mãe guardava todas essas coisas no seu coração.
\VS{52}Jesus crescia em sabedoria, em estatura, e em graça para com Deus e as pessoas.

\par }\Chap{3}{\PP \VerseOne{1}No ano quinze do império de Tibério César, sendo Pôncio Pilatos o governador da Judeia, Herodes tetraca da Galileia, e seu irmão Filipe o tetrarca da Itureia e da província de Traconites, e Lisânias o tetrarca de Abilene;
\VS{2}sendo Anás e Caifás os sumos sacerdotes, a palavra de Deus veio a João, filho de Zacarias, no deserto.
\VS{3}Ele percorreu toda a região ao redor do Jordão, pregando o batismo de arrependimento, para o perdão dos pecados;
\VS{4}conforme o que está escrito no livro das palavras do profeta Isaías, que diz: “Voz do que clama no deserto: Preparai o caminho do Senhor, endireitai as suas veredas.
\VS{5}Todo vale se encherá, e todo monte e colina se abaixará; os tortos serão endireitados, e os caminhos ásperos se suavizarão;
\VS{6}e todos verão
\FTNT{*}{{\FR 3:6 }
Lit. toda carne verá
} a salvação de Deus.”
\VS{7}{\ADD{João}} dizia às multidões que vinham ser batizadas por ele: “Raça de víboras, quem vos ensinou a fugir da ira que está para chegar?
\VS{8}Dai, pois, frutos condizentes com o arrependimento, e não comeceis a dizer em vós mesmos: ‘Temos Abraão por pai’; pois eu vos digo que até destas pedras Deus pode produzir
\FTNT{†}{{\FR 3:8 }
Lit. levantar
} filhos a Abraão.
\VS{9}E também o machado já está posto à raiz das árvores; portanto, toda árvore que não dá bom fruto é cortada e lançada ao fogo.
\VS{10}E as multidões lhe perguntavam: Então o que devemos fazer?
\VS{11}E ele lhes respondeu: Quem tiver duas túnicas, reparta com o não tem; e quem tiver comida, faça da mesma maneira.
\VS{12}E chegaram também uns publicanos, para serem batizados; e disseram-lhe: Mestre, que devemos fazer?
\VS{13}E ele lhes disse: Não exijais mais do que vos está ordenado.
\VS{14}E uns soldados também lhe perguntaram: E nós, que devemos fazer? E ele lhes disse: Não maltrateis nem defraudeis a ninguém; e contentai-vos com vossos salários.
\VS{15}E, enquanto o povo estava em expectativa, e todos pensando em seus corações se talvez João seria o Cristo,
\VS{16}João respondeu a todos: De fato, eu vos batizo com água, mas vem um que é mais poderoso do que eu, de quem eu não sou digno de desatar a tira das sandálias; ele vos batizará com o Espírito Santo e com fogo.
\VS{17}Sua pá está na sua mão, e limpará sua eira, e juntará o trigo no seu celeiro; porém queimará a palha com fogo que nunca se apaga.
\VS{18}Assim, exortando com muitas outras coisas,
{\ADD{João}} anunciava a boa notícia ao povo.
\VS{19}Porém, quando Herodes, o tetrarca, foi repreendido por ele por causa de Herodias, mulher do seu irmão Filipe, e por todas as maldades que Herodes havia feito,
\VS{20}acrescentou a todas ainda esta: encarcerou João na prisão.
\VS{21}E aconteceu que, quando todo o povo era batizado, Jesus
{\ADD{também}} foi batizado, e enquanto orava, o céu se abriu,
\VS{22}e o Espírito Santo desceu sobre ele em forma corpórea, como pomba; e veio do céu uma voz que dizia: Tu és o meu Filho amado, em ti me agrado.
\VS{23}E o mesmo Jesus começou quando tinha cerca de trinta anos, sendo (como se pensava) filho de José,
{\ADD{filho}} de Eli,
\VS{24}{\ADD{filho}} de Matate,
{\ADD{filho}} de Levi,
{\ADD{filho}} de Melqui,
{\ADD{filho}} de Janai,
{\ADD{filho}} de José.
\VS{25}{\ADD{filho}} de Matatias,
{\ADD{filho}} de Amós,
{\ADD{filho}} de Naum,
{\ADD{filho}} de Esli,
{\ADD{filho}} de Nagai.
\VS{26}{\ADD{filho}} de Maate,
{\ADD{filho}} de Matatias,
{\ADD{filho}} de Semei,
{\ADD{filho}} de José,
{\ADD{filho}} de Judá.
\VS{27}{\ADD{filho}} de Joaná,
{\ADD{filho}} de Resa,
{\ADD{filho}} de Zorobabel, e
{\ADD{filho}} de Salatiel,
{\ADD{filho}} de Neri.
\VS{28}{\ADD{filho}} de Melqui,
{\ADD{filho}} de Adi, e
{\ADD{filho}} de Cosã,
{\ADD{filho}} de Elmodã,
{\ADD{filho}} de Er.
\VS{29}{\ADD{filho}} de José,{\ADD{filho}} de Eliézer,
{\ADD{filho}} de Jorim,
{\ADD{filho}} de Matate,
{\ADD{filho}} de Levi.
\VS{30}{\ADD{filho}} de Simeão,
{\ADD{filho}} de Judá,
{\ADD{filho}} de José,
{\ADD{filho}} de Jonã,
{\ADD{filho}} de Eliaquim.
\VS{31}{\ADD{filho}} de Meleá,
{\ADD{filho}} de Mainã,
{\ADD{filho}} de Matatá,
{\ADD{filho}} de Natã,
{\ADD{filho}} de Davi.
\VS{32}{\ADD{Filho}} de Jessé,
{\ADD{filho}} de Obede,
{\ADD{filho}} de Boaz,
{\ADD{filho}} de Salmom,
{\ADD{filho}} de Naassom.
\VS{33}{\ADD{Filho}} de Aminadabe,
{\ADD{filho}} de Arão,{\ADD{filho}} de Esrom,
{\ADD{filho}} de Peres,
{\ADD{filho}} de Judá.
\VS{34}{\ADD{filho}} de Jacó,
{\ADD{filho}} de Isaque,
{\ADD{filho}} de Abraão,
{\ADD{filho}} de Terá,
{\ADD{filho}} de Naor.
\VS{35}{\ADD{filho}} de Serugue,
{\ADD{filho}} de Ragaú,
{\ADD{filho}} de Faleque,
{\ADD{filho}} de Éber,
{\ADD{filho}} de Salá.
\VS{36}{\ADD{Filho}} de Cainã,
{\ADD{filho}} de Arfaxade,
{\ADD{filho}} de Sem,
{\ADD{filho}} de Noé,
{\ADD{filho}} de Lameque.
\VS{37}{\ADD{Filho}} de Matusalém,
{\ADD{filho}} de Enoque,
{\ADD{filho}} de Jarede,
{\ADD{filho}} de Maleleel,
{\ADD{filho}} de Cainã.
\VS{38}{\ADD{filho}} de Enos,
{\ADD{filho}} de Sete,
{\ADD{filho}} de Adão,
{\ADD{e Adão}} de Deus.

\par }\Chap{4}{\PP \VerseOne{1}E Jesus, cheio do Espírito Santo, voltou do Jordão, e foi levado pelo Espírito ao deserto.
\VS{2}E por quarenta dias foi tentado pelo diabo; e não comeu coisa alguma naqueles dias; e terminados eles, depois teve fome.
\VS{3}Então o diabo lhe disse: Se tu és Filho de Deus, diz a esta pedra que se transforme em pão.
\VS{4}Jesus lhe respondeu: Está escrito que não só de pão viverá o ser humano, mas de toda palavra de Deus.
{\RQ{Deuteronômio 8:3
}}
\VS{5}E o diabo o levou a um alto monte, e lhe mostrou todos os reinos do mundo num instante.
\VS{6}E o diabo lhe disse: Darei a ti todo este poder e sua glória; pois a mim foram entregues, e os dou a quem quero.
\VS{7}Portanto, se tu me adorares, tudo será teu.
\VS{8}E Jesus lhe respondeu: Para trás de mim, Satanás!Porque está escrito: Adorarás ao Senhor teu Deus, e só a ele servirás.
{\RQ{Deuteronômio 6:13
}}
\VS{9}Então o levou a Jerusalém, e o pôs sobre a parte mais alta do templo, e disse-lhe: Se tu és o Filho de Deus, lança-te daqui abaixo.
\VS{10}Porque está escrito que: Mandará aos seus acerca de ti que te guardem.
{\RQ{Salmos 91:11
}}
\VS{11}E que: te sustentarão pelas mãos, para que nunca tropeces com o teu pé em alguma pedra.
{\RQ{Salmos 91:12
}}
\VS{12}Jesus lhe respondeu: Dito está: Não tentes
\FTNT{*}{{\FR 4:12 }
Ou: testes
} ao Senhor teu Deus.
{\RQ{Deuteronômio 6:16
}}
\VS{13}E quando o diabo acabou toda a tentação, ausentou-se dele por algum tempo.
\VS{14}Então Jesus, no poder do Espírito, voltou para a Galileia, e sua fama se espalhou por toda a região ao redor.
\VS{15}Ele ensinava em suas sinagogas, e era recebido com honra por todos.
\VS{16}Ele foi a Nazaré, onde havia sido criado, e entrou, conforme o seu costume, num dia de Sábado, na sinagoga; e levantou-se para ler.
\VS{17}Foi-lhe dado o livro do profeta Isaías; e quando abriu o livro, achou o lugar onde estava escrito:
\VS{18}O Espírito do Senhor
{\ADD{está}} sobre mim, porque ele me ungiu para evangelizar os pobres, me enviou para curar os contritos de coração;
\VS{19}para proclamar liberdade aos cativos e dar vista aos cegos, para pôr em liberdade os oprimidos, para proclamar o ano agradável do Senhor.
\FTNT{†}{{\FR 4:19 }
ano agradável do Senhor
Ou: ano do favor do Senhor
}
{\RQ{Isaías 61:1-2
}}
\VS{20}Ele fechou o livro, e devolveu-o ao assistente, e se sentou. Os olhos de todos na sinagoga estavam fixos nele.
\VS{21}Então ele começou a lhes dizer: Hoje esta escritura se cumpriu em vossos ouvidos.
\VS{22}E todos davam testemunho dele, e se admiravam das palavras graciosas que saíam da sua boca; e diziam: Não é este o filho de José?
\VS{23}E ele lhes disse: Sem dúvida me direis este provérbio: Médico, cura a ti mesmo; faz também aqui na tua terra natal todas as coisas que ouvimos terem sido feitas em Cafarnaum.
\VS{24}E disse: Em verdade vos digo que nenhum profeta é bem recebido na sua terra natal.
\VS{25}Porém em verdade vos digo que hava muitas viúvas em Israel nos dias de Elias, quando o céu se fechou por três anos e seis meses, de modo que em toda
{\ADD{aquela}} terra houve grande fome.
\VS{26}Mas a nenhuma delas Elias foi enviado, a não ser a uma mulher viúva de Sarepta de Sidom.
\VS{27}E havia muitos leprosos em Israel no tempo do profeta Eliseu; mas nenhum deles foi limpo, a não ser Naamã, o sírio.
\VS{28}E todos na sinagoga, quando ouviram essas coisas, encheram-se de ira.
\VS{29}Então se levantaram, expulsaram-no da cidade, e o levaram até o topo do monte em que sua cidade era construída, para dali o lançarem abaixo.
\VS{30}Porém ele passou por meio deles, e se retirou.
\VS{31}Ele desceu a Cafarnaum, cidade de Galileia; e
{\ADD{ali}} os ensinava nos sábados.
\VS{32}E admiravam o seu ensino, pois a sua palavra era com autoridade.
\VS{33}E estava na sinagoga um homem que tinha um espírito de um demônio imundo, e gritou com alta voz,
\VS{34}dizendo: Ah, que temos contigo, Jesus Nazareno? Vieste para nos destruir? Sei quem és: o Santo de Deus.
\VS{35}E Jesus o repreendeu, dizendo: Cala-te, e, sai dele. E o demônio o derrubou no meio
{\ADD{do povo}} , e saiu dele sem lhe fazer dano algum.
\VS{36}E espanto veio sobre todos; e falavam entre si uns aos outros, dizendo: Que palavra é esta, que até aos espíritos imundos ele manda com autoridade e poder, e eles saem?
\VS{37}E sua fama era divulgada em todos os lugares em redor daquela região.
\VS{38}{\ADD{Jesus}} levantou-se da sinagoga, e entrou na casa de Simão. A sogra de Simão estava doente de uma grande febre, e rogaram-lhe por ela.
\VS{39}Ele se inclinou a ela, e repreendeu a febre, que a deixou. Ela imediatamente se levantou e começou a lhes servir.
\VS{40}Quando o sol estava se pondo, troxeram-lhe todos os que estavam enfermos de varias doenças. Ele punha as mãos sobre cada um deles e os curava.
\VS{41}E também de muitos saíam demônios, gritando, e dizendo: Tu és o Cristo, o Filho de Deus. Ele os repreendia, e não os deixava falar, porque sabiam que ele era o Cristo.
\VS{42}E sendo já dia, ele saiu, e foi a um lugar deserto. As multidões o buscavam; então vieram a ele, e queriam detê-lo, para que não os deixasse.
\VS{43}Porém ele lhes disse: Também é necessário que eu anuncie a outras cidades o evangelho do Reino de Deus; pois para isso sou enviado.
\VS{44}E ele pregava nas sinagogas da Galileia.

\par }\Chap{5}{\PP \VerseOne{1}E aconteceu que as multidões o comprimiam para ouvir a palavra de Deus, quando ele estava junto ao lago de Genesaré.
\VS{2}Então ele viu dois barcos que estavam próximos
{\ADD{da margem}} do lago, e os pescadores haviam descido deles, e estavam lavando as redes.
\VS{3}Ele entrou num daqueles barcos, que era o de Simão, e lhe pediu que o afastasse um pouco da terra. Depois se sentou, e do barco ensinou às multidões.
\VS{4}Quando acabou de falar, disse a Simão: Vai à água profunda, e lançai as vossas redes para pescar.
\VS{5}Simão lhe respondeu: Mestre, trabalhamos toda a noite, mas nada pegamos; mas pela tua palavra lançarei a rede.
\VS{6}Eles fizeram assim, e recolheram grande quantidade de peixes, e sua rede estava se rompendo.
\VS{7}Então acenaram aos companheiros que estavam no outro barco, para que viessem os ajudar. Vieram, e encheram ambos os barcos, de tal maneira que quase afundavam.
\VS{8}Quando Simão Pedro viu{\ADD{isso}} , prostrou-se aos pés de Jesus, dizendo: Afasta-te de mim, Senhor, porque sou um homem pecador.
\VS{9}Pois ele estava cheio de espanto, como todos os que estavam com ele, por causa da coleta de peixes que haviam feito.
\VS{10}E semelhantemente, também, Tiago e João, filhos de Zebedeu, que eram colegas de Simão. E Jesus disse a Simão: Não temas, de agora em diante pescarás gente.
\VS{11}Eles levaram os barcos para a terra, deixaram tudo, e o seguiram.
\VS{12}E aconteceu que, quando ele estava numa daquelas cidades, eis que um homem cheio de lepra viu Jesus; então prostrou-se sobre o rosto, e lhe rogou, dizendo: Senhor, se quiseres, podes me tornar limpo.
\VS{13}Ele estendeu a mão, e o tocou, dizendo: Quero; sê limpo. E logo a lepra desapareceu dele.
\VS{14}E mandou-lhe que não o dissesse a ninguém: Mas vai, mostra-te ao sacerdote, e oferece por tua purificação, como Moisés mandou, para lhes dar testemunho.
\VS{15}Porém a sua fama se espalhava ainda mais; e ajuntavam-se muitas multidões para o ouvir, e para serem por ele curadas de suas enfermidades.
\VS{16}Porém ele se retirava para os desertos, e
{\ADD{ali}} orava.
\VS{17}E aconteceu num daqueles dias que estava ensinando, e estavam
{\ADD{ali}} sentados fariseus e instrutores da lei, que tinham vindo de todas as aldeias da Galileia, da Judeia, e de Jerusalém; e o poder do Senhor estava
{\ADD{ali}} para os curar.
\VS{18}E eis que uns homens traziam numa cama um homem que estava paralítico; e procuravam levá-lo para dentro, a fim de o porem diante dele.
\VS{19}Como não acharam por onde pudessem levá-lo para dentro, por causa da multidão, subiram ao telhado, e pelas telhas o baixaram com o leito para o meio
{\ADD{do povo}} , diante de Jesus.
\VS{20}Ele viu a fé deles, e disse-lhe: Homem, teus pecados são perdoados.
\VS{21}E os escribas e os fariseus começaram a questionar, dizendo: Quem é este que fala blasfêmias? Quem pode perdoar pecados, senão só Deus?
\VS{22}Porém Jesus, conhecendo os seus pensamentos, respondeu-lhes: O que pensais em vossos corações?
\VS{23}O que é mais fácil? Dizer: Teus pecados são perdoados? Ou dizer: Levanta-te, e anda?
\VS{24}Para que saibais que o Filho do homem tem poder sobre a terra para perdoar pecados, (disse ao paralítico:) Eu te digo, levanta-te, toma o teu leito, e vai para a tua casa.
\VS{25}E ele levantou-se imediatamente diante deles, tomou o
{\ADD{leito}} em que estava deitado, e foi para a sua casa, glorificando a Deus.
\VS{26}E todos ficaram admirados, e glorificaram a Deus; e ficaram cheios de temor, dizendo: Hoje vimos maravilhas!
\VS{27}Depois destas coisas, ele saiu; e viu um publicano, chamado Levi, sentado na coletoria de tributos, e disse-lhe: Segue-me.
\VS{28}Ele deixou tudo, levantou-se, e o seguiu.
\VS{29}E Levi lhe fez um grande banquete em sua casa; e havia
{\ADD{ali}} uma grande multidão de publicanos e de outros que estavam sentados com eles
{\ADD{à mesa}} .
\VS{30}E seus escribas e fariseus murmuravam contra os seus discípulos, dizendo: Por que comeis e bebeis com publicanos e pecadores?
\VS{31}Jesus lhes respondeu: Os que estão sãos não precisam de médico, mas sim os que estão doentes.
\VS{32}Eu não vim para chamar os justos, mas sim, os pecadores, ao arrependimento.
\VS{33}Então eles lhe disseram: Por que os discípulos de João jejuam muitas vezes, e fazem orações, como também os dos fariseus, mas os teus comem e bebem?
\VS{34}Mas ele lhes disse: Podeis vós fazer jejuar os convidados do casamento,
\FTNT{*}{{\FR 5:34 }
convidados do casamento
Lit. filhos da câmara nupcial
} enquanto o noivo está com eles?
\VS{35}Porém virão dias em que o esposo lhes será tirado; então naqueles dias jejuarão.
\VS{36}E dizia-lhes também uma parábola: Ninguém põe remendo de pano novo em roupa velha; de outra maneira, o novo rompe, e o remendo do novo não será adequado ao velho.
\VS{37}E ninguém põe vinho novo em odres velhos; de outra maneira, o vinho novo romperá os odres, e se derramará, e os odres são destruídos.
\VS{38}Mas o vinho novo deve ser posto em odres novos, e ambos se conservam.
\VS{39}E ninguém que beber do velho quer logo o novo; porque diz: O velho é melhor.

\par }\Chap{6}{\PP \VerseOne{1}E aconteceu que, no segundo sábadodepois do primeiro,
{\ADD{Jesus}} passou pelas plantações, e seus discípulos iam arrancando espigas, e comiam, debulhando-as com as mãos.
\VS{2}E alguns dos fariseus lhes disseram: Por que fazeis o que não é lícito fazer nos sábados?
\VS{3}E Jesus lhes respondeu: Nunca lestes isto, o que Davi fez quando teve fome, ele e os que com ele estavam?
\VS{4}Como entrou na casa de Deus, tomou, e comeu os pães da apresentação, e deu também aos que estavam com ele,
{\ADD{pães}} que não é lícito comer, a não ser só os sacerdotes?
\VS{5}E dizia-lhes: O Filho do homem é Senhor até do sábado.
\VS{6}E aconteceu também em outro sábado que entrou na sinagoga, e estava ensinado; e ali estava um homem que tinha a mão direita definhada.
\VS{7}E os escribas e fariseus o observavam, se o curaria no sábado; para acharem de que o acusar.
\VS{8}Mas ele bem sabia dos seus pensamentos; e disse ao homem que tinha a mão definhada: Levanta-te, e põe-te em pé no meio. E ele se levantou e se pôs de pé.
\VS{9}Então Jesus lhes disse: Uma coisa vos perguntarei: É lícito nos sábados fazer bem, ou fazer mal? Salvar uma pessoa, ou matá-la?
\VS{10}E ele, olhando para todos em redor, disse ao homem: Estende a tua mão. E ele assim o fez; e a mão foi lhe restituída sã como a outra.
\VS{11}Então encheram-se de ira; e conversaram uns com os outros
{\ADD{sobre}} o que fariam a Jesus.
\VS{12}E aconteceu que naqueles dias ele foi ao monte para orar; e passou a noite orando a Deus.
\VS{13}E quando já era dia, chamou a si os seus discípulos, e escolheu doze deles, a quem também chamou de apóstolos:
\VS{14}Simão, a quem também chamou de Pedro, e seu irmão André; Tiago, e João; Filipe, e Bartolomeu.
\VS{15}Mateus, Tomé; Tiago
{\ADD{filho}} de Alfeu, e Simão chamado “o Zelote”.
\VS{16}Judas de Tiago, e Judas Iscariotes, o mesmo que foi o traidor.
\VS{17}Ele desceu com eles, parou num lugar plano, e também
{\ADD{com ele}} uma grande quantidade dos seus discípulos, e grande multidão do povo de toda a Judeia, de Jerusalém, da costa marítima de Tiro, e de Sidom,
\VS{18}que tinham vindo para o ouvir, e para serem curados das suas enfermidades, como também os atormentados por espíritos imundos; e foram curados.
\VS{19}E toda a multidão procurava tocá-lo; porque dele saía poder, e curava a todos.
\VS{20}Ele levantou os olhos aos seus discípulos, e disse: Benditos
\FTNT{*}{{\FR 6:20 }
Ou: Bem-aventurados; também nos versos seguintes
} sois vós, os pobres, porque o Reino de Deus é vosso.
\VS{21}Benditos sois vós que agora tendes fome, porque sereis saciados. Benditos sois vós que agora chorais, porque rireis.
\VS{22}Benditos sereis quando as pessoas vos odiarem, quando vos separarem, vos insultarem, e rejeitarem o vosso nome como mau, por causa do Filho do homem.
\VS{23}Contentai-vos nesse dia, e saltai de alegria, porque eis que grande é a vossa recompensa nos céus; pois assim os pais deles faziam aos profetas.
\VS{24}Mas ai de vós, ricos, porque já tendes a vossa consolação.
\VS{25}Ai de vós que estais saciados, porque tereis fome. Ai de vós que agora rides, porque lamentareis, e chorareis.
\VS{26}Ai de vós quando todas as pessoas falarem bem de vós; porque assim os pais deles faziam aos falsos profetas.
\VS{27}Mas a vós, que estais ouvindo, digo: amai os vossos inimigos; fazei bem aos que vos odeiam;
\VS{28}bendizei aos que vos maldizem, e orai pelos que vos maltratam.
\VS{29}Ao que te ferir numa face, oferece-lhe também a outra; e ao que te tirar a capa, não recuses a túnica.
\VS{30}Dá a quem te pedir; e ao que te tomar o
{\ADD{que é}} teu, não o peças de volta.
\VS{31}E como vós quereis que as pessoas vos façam, fazei-lhes vós também da mesma maneira.
\VS{32}E se amardes aos que vos amam, que mérito tereis? Pois também os pecadores amam os que os amam.
\VS{33}E se fizerdes bem aos que vos fazem bem, que mérito tereis? Pois também os pecadores fazem o mesmo.
\VS{34}E se emprestardes àqueles de quem esperais receber de volta, que mérito tereis? Pois também os pecadores emprestam aos pecadores, para receberam de volta o tanto equivalente.
\VS{35}Em vez disso, amai aos vossos inimigos, fazei o bem, e emprestai, sem nada esperar disso; e grande será a vossa recompensa, e sereis filhos do Altíssimo; porque é benigno
{\ADD{até}} para com os ingratos e maus.
\VS{36}Portanto, sede misericordiosos, como também o vosso Pai é misericordioso.
\VS{37}Não julgueis, e não sereis julgados; não condeneis, e não sereis condenados; liberai, e vos liberarão.
\VS{38}Dai, e será vos dado; medida boa, comprimida, sacudida e transbordante vos darão no vosso colo; pois com a mesma medida que medirdes vos medirão de volta.
\VS{39}E disse-lhesuma parábola: Acaso pode o cego guiar
{\ADD{outro}} cego? Não cairão ambos no buraco?
\VS{40}O discípulo não está acima do seu mestre; mas todo aquele que estiver completamente capacitado será como o seu mestre.
\VS{41}E por que tu prestas atenção no cisco que está no olho do teu irmão, e não enxergas a trave que está no teu próprio olho?
\VS{42}Ou como podes dizer a teu irmão: “Irmão, deixa-me tirar o cisco que está no teu olho”, se tu mesmo não prestas atenção na trave que está no teu olho? Hipócrita, tira primeiro a trave de teu olho, e então verás bem para tirar o cisco que está no olho do teu irmão.
\VS{43}Pois não há boa árvore que dê mau fruto, nem árvore má que dê bom fruto.
\VS{44}Porque cada árvore se conhece pelo seu próprio fruto; pois não se colhem figos dos espinheiros, nem se tiram uvas dos cardos.
\VS{45}A boa pessoa tira o bem do bom tesouro do seu coração, e a pessoa má tira o mal do mau tesouro do seu coração; pois a sua boca fala daquilo que o coração tem em abundância.
\VS{46}E por que me chamais: “Senhor!”, “Senhor!”, e não fazeis o que eu digo?
\VS{47}Todo aquele que vem a mim, ouve as minhas palavras, e as pratica, eu vos mostrarei a quem é semelhante:
\VS{48}Semelhante é ao homem que construiu uma casa, cavou, abriu bem fundo, e pôs o fundamento sobre a rocha; e quando veio a enchente, a corrente bateu com ímpeto naquela casa, e não a pôde abalar, porque estava fundada sobre a rocha.
\VS{49}Mas aquele que ouve e não pratica é semelhante ao homem que construiu uma casa sobre a terra, sem fundamento, na qual a corrente bateu com ímpeto, e logo caiu; e foi grande a queda daquela casa.

\par }\Chap{7}{\PP \VerseOne{1}E depois de acabar todos os seus discursos aos ouvidos do povo, ele entrou em Cafarnaum.
\VS{2}E o servo de um certo centurião, a quem muito estimava, que estava doente, e a ponto de morrer.
\VS{3}E quando
{\ADD{o centurião}} ouviu falar sobre Jesus, enviou-lhe uns anciãos dos judeus, rogando-lhe que viesse curar o seu servo.
\VS{4}Eles vieram a Jesus, e rogaram-lhe com urgência, dizendo: Ele é digno de lhe concederes isto,
\VS{5}porque ele ama a nossa nação, e ele mesmo construiu a nossa sinagoga.
\VS{6}Jesus foi com eles; mas quando já não estava longe da casa, o centurião enviou-lhe uns amigos, dizendo-lhe: “Senhor, não te incomodes, porque não sou digno que entres abaixo do meu telhado.
\VS{7}Por isso que nem mesmo me considerei digno de vir a ti; mas diz uma palavra, e o meu servo sarará.
\VS{8}Porque também eu sou homem subordinado à autoridade, e tenho soldados sob o meu comando, e digo a este: 'Vai', e ele vai; e a outro: 'Vem', e ele vem; e a meu servo: 'Faz isto', e ele faz”.
\VS{9}Quando Jesus ouviu isso, admirou-se dele; então se virou, e disse à multidão que o seguia: Digo-vos que nem mesmo em Israel achei tanta fé.
\VS{10}E quando os que foram enviados voltaram para casa, acharam com boa saúde o servo antes doente.
\VS{11}E aconteceu no dia seguinte, que
{\ADD{Jesus}} ia a uma cidade chamada Naim, e com ele iam muitos de seus discípulos e uma grande multidão.
\VS{12}E quando chegou perto da porta da cidade, eis que levavam um defunto, filho único de sua mãe, que
{\ADD{era}} viúva; e com ela
{\ADD{ia}} grande multidão da cidade.
\VS{13}Quando o Senhor a viu, comoveu-se de intima compaixão por ela, e disse-lhe: Não chores.
\VS{14}Então se aproximou, tocou o caixão (e os que a levavam, pararam), e disse: Jovem, a ti eu digo: levanta-te.
\VS{15}E o defunto se sentou e começou a falar; e ele o entregou à sua mãe.
\VS{16}Todos se encheram de temor, e glorificavam a Deus, dizendo: “Um grande profeta se levantou entre nós, e Deus visitou o seu povo!”
\VS{17}E essa sua fama correu por toda a Judeia, e por toda a região em redor.
\VS{18}E os discípulos de João lhe anunciaram essas coisas.
\VS{19}Então João chamou dois de seus discípulos, e os enviou a Jesus, dizendo: “És tu aquele que havia de vir, ou esperamos outro?”
\VS{20}E quando aqueles homens vieram a ele, disseram: “João Batista nos enviou a ti, dizendo: 'És tu aquele que havia de vir, ou esperamos outro?'”
\VS{21}E naquela mesma hora ele curou a muitos de enfermidades, males, e espíritos maus, e a deu vista a muitos cegos.
\VS{22}E Jesus lhes respondeu: Ide, e anunciai a João as coisas que tendes visto e ouvido: que os cegos veem, os mancos andam, os leprosos são purificados, os surdos ouvem, os mortos ressuscitam, e aos pobres é anunciado o Evangelho.
\VS{23}E bendito
\FTNT{*}{{\FR 7:23 }
Ou: bem-aventurado
} é aquele que não se ofender em mim.
\VS{24}E quando os mensageiros de João se foram,
{\ADD{Jesus}} começou a dizer às multidões acerca de João: Que saístes para ver no deserto? Alguma cana sacudida pelo vento?
\VS{25}Mas que saístes para ver? Um homem vestido de roupas delicadas? Eis que os que
{\ADD{vestem}} roupas delicadas e vivem no luxo estão nos palácios reais.
\VS{26}Mas que saístes para ver? Um profeta? Sim, vos digo, e muito mais que profeta.
\VS{27}Este é aquele de quem está escrito: Eis que envio o meu mensageiro adiante de tua face, o qual preparará o teu caminho diante de ti.
\VS{28}Pois eu vos digo, que dentre os nascidos de mulheres, não há profeta maior que João, o Batista; mas o menor no reino dos céus é maior que ele.
\VS{29}E todo o povo e os publicanos ouviram, concordaram que Deus era justo, e foram batizados com o batismo de João.
\VS{30}Mas os fariseus e os estudiosos da Lei rejeitaram o conselho de Deus contra si mesmos, e não foram batizados por ele.
\VS{31}E o Senhor disse: “A quem, pois, compararei as pessoas desta geração? E a quem são semelhantes?
\VS{32}São semelhantes às crianças sentadas na praça, que gritam umas às outras: 'Nós vos tocamos flautas, mas não dançastes; nós vos cantamos lamentações, mas não chorastes'.
\VS{33}Porque veio João Batista, que não comia pão, nem bebia vinho, e dizeis: 'Ele tem demônio';
\VS{34}Veio o Filho do homem, que come e bebe, e dizeis: 'Eis um homem comilão e bebedor de vinho, amigo dos publicanos e dos pecadores'.
\VS{35}Mas a sabedoria foi considerada justa por todos os seus filhos”.
\VS{36}E um dos fariseus lhe pediu que comesse com ele; então entrou na casa do fariseu, e sentou-se
{\ADD{à mesa}} .
\VS{37}E eis que uma mulher da cidade, que era pecadora, quando soube que ele estava na casa do fariseu, trouxe um vaso de alabastro de óleo perfumado.
\VS{38}E estando atrás, aos seus pés, chorando, começou a molhar-lhe os pés com lágrimas; e os enxugava com os cabelos da sua cabeça; e beijava-lhe os pés, e os ungia com o óleo perfumado.
\VS{39}E quando o fariseu que o havia convidado viu
{\ADD{isso}} , falou consigo mesmo: Se ele fosse profeta, saberia quem e qual é a mulher que o toca; porque ela é pecadora.
\VS{40}E Jesus lhe respondeu: “Simão, tenho uma coisa a te dizer”; e ele disse: “Dize-a, Mestre”.
\VS{41}{\ADD{Jesus}} disse: “Certo credor tinha dois devedores. Um lhe devia quinhentas moedas de prata, e o outro cinquenta.
\VS{42}Como não tinham com que pagar, perdoou-lhes
{\ADD{a dívida}} de ambos. Dize, pois, qual deles o amará mais?”
\VS{43}Simão respondeu: Tenho para mim que aquele a quem mais perdoou. E ele lhe disse: Julgaste corretamente.
\VS{44}Ele se voltou à mulher, e disse a Simão: Vês tu esta mulher? Entrei na tua casa, e não me deste água para os pés; porém esta molhou os meus pés com lágrimas, e os enxugou com os seus cabelos da cabeça.
\VS{45}Tu não me beijaste; porém esta, desde que entrou, não parou de me beijar os pés.
\VS{46}Não ungiste a minha cabeça com óleo, porém esta ungiu os meus pés com óleo perfumado.
\VS{47}Por isso te digo, os muitos pecados dela são perdoados, porque muito amou; mas ao que pouco se perdoa, pouco ama.
\VS{48}E disse a ela: Os teus pecados são perdoados.
\VS{49}E os que estavam sentados começaram a dizer entre si: Quem é este, que até perdoa pecados?
\VS{50}E disse à mulher: A tua fé te salvou; vai em paz.

\par }\Chap{8}{\PP \VerseOne{1}E aconteceu depois disso, que
{\ADD{Jesus}} andava de cidade em cidade, e de aldeia em aldeia, pregando e anunciando o evangelho do reino de Deus; e com ele os doze,
\VS{2}e algumas mulheres que haviam sido curadas de espíritos malignos e de enfermidades: Maria, chamada Madalena, da qual saíram sete demônios;
\VS{3}e Joana, a mulher de Cuza, mordomo de Herodes; e Susana, e muitas outras, que o serviam com seus bens.
\VS{4}E quando se ajuntou uma multidão, e vindo a ele de todas as cidades, disse por parábola:
\VS{5}Saiu um semeador a semear sua semente; e quando semeava, caiu uma
{\ADD{parte}} junto ao caminho, e foi pisada, e as aves do céu a comeram.
\VS{6}E outra
{\ADD{parte}} caiu sobre a pedra; e nascida, secou-se, porque não tinha umidade.
\VS{7}E outra
{\ADD{parte}} caiu entre espinhos, e nascendo com ela, os espinhos a sufocaram.
\VS{8}E outra
{\ADD{parte}} caiu em boa terra, e nascida, produziu fruto a cem por um. Depois que disse coisas, exclamava: Quem tem ouvidos para ouvir, ouça.
\VS{9}E os seus discípulos lhe perguntaram: Que parábola é esta?
\VS{10}Ele disse: A vós é dado conhecer os mistérios do reino de Deus; mas aos outros por parábolas, para que vendo, não vejam; e ouvindo, não entendam.
\VS{11}Esta é, pois, a parábola: a semente é a palavra de Deus.
\VS{12}Os que estão junto ao caminho são os que ouvem; depois vem o diabo, e tira-lhes do coração a palavra, para que não creiam nem se salvem.
\VS{13}E os que estão sobre a pedra são os que, quando ouvem, recebem a palavra com alegria, mas esses não têm raiz, pois creem por algum tempo, mas no tempo da tentação se desviam.
\VS{14}E o que caiu entre espinhos são os que ouviram, e indo, sufocam-se com as preocupações, e riquezas, e prazeres da vida, e não produzem bom fruto.
\VS{15}E o que caiu em boa terra, estes são os que, ouvindo a palavra, conservam-na num coração honesto e bom, e dão fruto que permanece.
\VS{16}E ninguém, quando acende uma lâmpada, cobre-a com algum vaso, ou a põe debaixo da cama; em vez disso põe na luminária, para que os que entram vejam a luz.
\VS{17}Porque não há coisa oculta que não venha a se manifestar; nem coisa escondida que não venha a ser conhecida, e chegada à luz.
\VS{18}Portanto, prestai atenção a como ouvis; pois àquele que tiver lhe será dado; e àquele que não tiver, até o que lhe parece ter, dele será tirado.
\VS{19}E vieram a ele sua mãe e seus irmãos, mas não conseguiam chegar a ele por causa da multidão.
\VS{20}E foi-lhe anunciado, dizendo: “Tua mãe e teus irmãos estão fora, e querem te ver”.
\VS{21}Porém, ele lhes respondeu: “Minha mãe e meus irmãos são aqueles que ouvem a palavra de Deus e a praticam”.
\VS{22}E aconteceu, num daqueles dias, que
{\ADD{Jesus}} entrou em um barco com seus discípulos, e disse-lhes: “Passemos para o outro lado do lago”. E partiram.
\VS{23}E enquanto navegavam,
{\ADD{Jesus}} adormeceu; então desceu uma tempestade de vento no lago, e enchiam-se
{\ADD{de água}} , e estavam em perigo.
\VS{24}Então se aproximaram dele e o despertaram, dizendo: “Mestre, Mestre, estamos perecendo!” E ele se levantou, repreendeu o vento e as ondas da água; e pararam, e fez-se bonança.
\VS{25}E disse-lhes: “Onde está a vossa fé?” Mas eles, temendo, admiravam-se, e diziam uns aos outros: “Quem é este, que manda até aos ventos, e à água, e lhe obedecem?”
\VS{26}E navegaram para a terra dos gadarenos, que é do lado oposto à Galileia.
\VS{27}E quando ele desceu à terra, veio da cidade ao seu encontro um homem, que havia muito tempo que tinha demônios, e não andava vestido, nem habitava em casa nenhuma, mas sim, nos sepulcros.
\VS{28}Ele viu Jesus, então exclamou, prostrou-se diante dele, e disse com grande voz: “O que eu tenho contigo, Jesus, Filho do Deus Altíssimo? Rogo-te que não me atormentes.”
\VS{29}Porque ele havia mandado ao espírito imundo que saísse daquele homem; pois muitas vezes apoderava-se dele. E o mantinham preso com cadeias e grilhões; porém ele quebrava o que lhe prendia, e era conduzido pelo demônio aos desertos.
\VS{30}E Jesus lhe perguntou: “Qual é teu nome?” E ele disse: “Legião” (porque muitos demônios tinham entrado nele).
\VS{31}E rogavam-lhe que não os mandasse ir para o abismo.
\VS{32}E havia ali uma manada de muitos porcos, que pastavam no monte; e rogaram-lhe que lhes concedesse entrar neles; e ele lhes concedeu.
\VS{33}Os demônios saíram daquele homem e entraram nos porcos; e a manada se lançou de um precipício no lago e se afogou.
\VS{34}Quando os que cuidavam
{\ADD{dos porcos}} viram o que havia acontecido, fugiram; então foram anunciar na cidade e nos campos.
\VS{35}E saíram para ver o acontecido, e vieram até Jesus; e acharam ao homem de quem os demônios tinham saído, vestido e em pleno juízo, sentado aos pés de Jesus; e temeram.
\VS{36}E os que tinham visto contaram-lhes também como aquele endemoninhado havia sido liberto.
\VS{37}E toda o povo da terra dos gadarenos ao redor lhe rogou que se retirasse deles; porque foram tomados por grande temor. Então ele entrou no barco e voltou.
\VS{38}E aquele homem, de quem os demônios tinham saído, rogou-lhe para que lhe permitisse estar com ele; mas Jesus o despediu, dizendo:
\VS{39}“Volta para tua casa, e conta como foram grandes as coisas que Deus te fez”. Ele se foi, e pregou por toda a cidade, como foram grandes as coisas que Jesus lhe tinha feito.
\VS{40}E aconteceu que, quando Jesus voltou, a multidão o recebeu; porque todos o estavam esperando.
\VS{41}E eis que chegou um homem, cujo nome era Jairo, e era chefe da sinagoga. Então ele se prostrou aos pés de Jesus e rogou-lhe que entrasse em sua casa;
\VS{42}porque tinha uma filha única, com doze anos de idade, que estava a ponto de morrer. E enquanto ele ia, as multidões o apertavam.
\VS{43}E uma mulher que tinha um fluxo de sangue, havia doze anos, e já tinha gastado todos os seus pertences com médicos, e não pôde ser curada por nenhum.
\VS{44}Ela se aproximou
{\ADD{de Jesus}} por detrás, e tocou a borda de sua roupa; e logo o fluxo de seu sangue parou.
\VS{45}E Jesus disse: “Quem me tocou?” Enquanto todos negavam, disse Pedro e os que com ele estavam: Mestre, as multidões te apertam e empurram, e dizes: Quem me tocou?
\VS{46}Jesus disse: “Alguém me tocou, pois sei que poder saiu de mim”.
\VS{47}Então a mulher, vendo que não podia se esconder, veio tremendo, prostrou-se diante dele, e declarou-lhe diante de todo o povo a causa por que o havia tocado, e como logo havia sarado.
\VS{48}E ele lhe disse: “Tem bom ânimo, filha; a tua fé te sarou; vai em paz”.
\VS{49}Estando ele ainda falando, veio um da
{\ADD{casa}} do chefe da sinagoga, que lhe disse: “Tua filha já está morta, não incomodes o Mestre.”
\VS{50}Porém quando Jesus o ouviu, respondeu-lhe: “Não temas; crê somente, e ela será curada”.
\VS{51}E quando entrou na casa, a ninguém deixou entrar, a não ser a Pedro, Tiago, João, e ao pai e à mãe da menina.
\VS{52}E todos choravam e lamentavam por ela; mas ele disse: “Não choreis; ela não está morta, mas dorme”.
\VS{53}E riam dele, porque sabiam que estava morta.
\VS{54}Porém ele, depois de expulsá-los todos, segurou a mão dela e clamou: “Levanta-te, menina”.
\VS{55}Então seu espírito voltou, e ela logo se levantou; e
{\ADD{Jesus}} mandou que dessem de comer a ela.
\VS{56}E seus pais ficaram maravilhados, mas ele lhes mandou que a ninguém dissessem o que havia acontecido.

\par }\Chap{9}{\PP \VerseOne{1}{\ADD{Jesus}} convocou seus doze discípulos, e lhes deu poder e autoridade sobre todos os demônios, e para curarem enfermidades.
\VS{2}E os enviou para pregar o reino de Deus e para curar os enfermos.
\VS{3}E disse-lhes: “Não tomeis nada convosco para o caminho: nem vara, nem bolsa, nem pão, nem dinheiro, nem tenhais duas túnicas cada um.
\VS{4}E em qualquer casa que entrardes, ficai ali, e dali saí.
\VS{5}E a todo os que não vos receberem, quando sairdes daquela cidade, sacudi até o pó dos vossos pés, em testemunho contra eles”.
\VS{6}Então eles partiram, e percorreram todas as aldeias, anunciando o Evangelho, e curando em todos os lugares.
\VS{7}E o tetrarca Herodes ouvia falar todas as coisas que ele fazia; e estava perplexo, porque alguns diziam que João tinha ressuscitado dos mortos;
\VS{8}e outros, que Elias havia aparecido; e outros, que algum profeta dos antigos havia ressuscitado.
\VS{9}E Herodes disse: “A João mandei degolar; quem, pois, é este, de quem tais coisas ouço?” E procurava vê-lo.
\VS{10}E quando os apóstolos, voltaram, contaram
{\ADD{a Jesus}} todas as coisas que tinham feito. Então ele os tomou consigo, e retirou-se à parte a um lugar deserto de uma cidade chamada Betsaida.
\VS{11}E quando as multidões souberam, seguiram-no. Ele as recebeu, e lhes falava do Reino de Deus, e curava aos que precisavam de cura.
\VS{12}E o dia já começava a declinar. Então os doze se aproximaram dele, e lhe disseram: “Despede a multidão, para irem aos lugares e aldeias ao redor, se agasalhem, e achem o que comer; porque aqui estamos em lugar deserto”.
\VS{13}Mas ele lhes disse: “Dai-lhes vós de comer”. E eles disseram: “Não temos mais que cinco pães e dois peixes; a não ser se formos nós mesmos comprar de comer para todo este povo”.
\VS{14}Porque havia ali quase cinco mil homens. Então disse aos seus discípulos: “Fazei-os se sentarem em grupos de cinquenta em cinquenta”.
\VS{15}E assim procederam, fazendo todos se sentarem.
\VS{16}Então ele tomou os cinco pães e os dois peixes e, olhando para o céu, abençoou-os e partiu-os, e os deu a seus discípulos, para
{\ADD{os}} porem diante da multidão.
\VS{17}E todos comeram, e saciaram-se; e levantaram doze cestos de pedaços que sobraram.
\VS{18}E aconteceu que, quando ele estava orando só, e seus discípulos estavam com ele, que ele lhes perguntou: “Quem as multidões dizem que eu sou?”
\VS{19}E eles responderam: “{\ADD{Uns}} , João Batista; outros, Elias; e outros, que algum dos profetas antigos ressuscitou”.
\VS{20}E disse-lhes: “E vós, quem dizeis que eu sou?” E Pedro respondeu: “O Cristo de Deus”.
\VS{21}Então ele os alertou, e mandou-lhes que não dissessem a ninguém,
\VS{22}dizendo: “É necessário que o Filho do homem sofra muitas coisas; que seja rejeitado pelos anciãos, pelos chefes dos sacerdotes, e pelos escribas; que seja morto, e que seja ressuscitado ao terceiro dia”.
\VS{23}E dizia a todos: “Se alguém quer vir após mim, negue-se a si mesmo, tome cada dia sua cruz, e siga- me.
\VS{24}Pois quem quiser salvar sua vida a perderá; porém quem, por minha causa, perder a sua vida, esse a salvará.
\VS{25}Porque, que aproveita para alguém ganhar o mundo todo, e perder ou prejudicar a si mesmo?
\VS{26}Pois quem se envergonhar de mim e das minhas palavras, dele o Filho do homem se envergonhará, quando vier na sua glória, e na do Pai, e dos santos anjos.
\VS{27}E em verdade vos digo que, dentre os que aqui estão, há alguns que não experimentarão a morte, até que vejam o Reino de Deus”.
\VS{28}E aconteceu que, quase oito dias depois dessas palavras, ele tomou consigo a Pedro, João, e a Jacó, e subiu ao monte para orar.
\VS{29}E ele, enquanto estava orando, a aparência do seu rosto se transfigurou, e sua roupa
{\ADD{ficou}} branca e brilhante.
\VS{30}E eis que estavam falando com ele dois homens, que eram Moisés e Elias,
\VS{31}os quais apareceram com glória, e falavam de sua partida, que estava para se cumprir em Jerusalém.
\VS{32}Pedro e os que estavam com ele estavam cheios de sono; e quando despertaram, viram a glória dele, e aqueles dois homens que estavam com ele.
\VS{33}E aconteceu que, quando eles estavam saíndo da presença dele, Pedro disse a Jesus: “Mestre, é bom estarmos nós aqui. Façamos três tendas, uma para ti, uma para Moisés, e uma para Elias”, não sabendo o que dizia.
\VS{34}E enquanto ele dizia isso, veio uma nuvem que os cobriu com a sua sombra; e, quando entraram na nuvem, temeram.
\VS{35}E veio uma voz da nuvem, que dizia: “Este é o meu Filho amado; a ele ouvi”.
\VS{36}E, depois de terminada aquela voz, Jesus se encontrava só; e eles se calaram, e por aqueles dias não contaram a ninguém nada do que tinham visto.
\VS{37}E aconteceu no dia seguinte que, descendo eles do monte, uma grande multidão lhe saiu ao encontro.
\VS{38}E eis que um homem da multidão clamou: “Mestre, rogo-te que vejas a meu filho, que é o único que tenho,
\VS{39}e eis que um espírito o toma, e de repente grita, e o convulsiona até espumar, e somente sai dele quando o maltrata severamente.
\VS{40}E roguei aos teus discípulos que o exupulsassem, mas não conseguiram”.
\VS{41}E Jesus respondeu: “Ó geração incrédula e perversa! Até quando estarei ainda convosco, e vos suportarei? Traze aqui teu filho”.
\VS{42}E quando ainda vinha chegando, o demônio o derrubou, e o convulsionou. Porém Jesus repreendeu o espírito imundo, curou ao menino, e o devolveu ao seu pai.
\VS{43}E todos ficaram perplexos com a grandeza de Deus. E enquanto todos se maravilhavam de tudo o que Jesus fazia, disse aos seus discípulos:
\VS{44}“Ponde vós em vossos ouvidos estas palavras: que o Filho do homem será entregue nas mãos dos homens”.
\VS{45}Mas eles não entendiam essa palavra, e era-lhes encoberta, para que não a compreendessem; e temiam perguntar-lhe sobre essa palavra.
\VS{46}E levantou-se entre eles uma discussão sobre qual deles seria o maior.
\VS{47}Mas Jesus, conhecendo o pensamento de seus corações, pegou uma criança, e a pôs junto a si.
\VS{48}Então disse-lhes: “Quem receber esta criança em meu nome, a mim me recebe; e quem receber a mim, recebe o que me enviou; porque aquele que entre todos vós for o menor, esse será grande”.
\VS{49}João respondeu: “Mestre, vimos um que em teu nome expulsava os demônios, e nós o proibimos, porque não segue conosco”.
\VS{50}E Jesus lhe disse: “Não o proibais, porque quem não é contra nós, é por nós”.
\VS{51}E aconteceu que, quando completaram-se os dias em que ele viria a ser recebido no alto, ele se determinou a ir para Jerusalém.
\VS{52}E mandou mensageiros adiante de si;
\FTNT{*}{{\FR 9:52 }
Lit. adiante de sua face
} e eles se foram, e entraram numa aldeia de samaritanos, para lhe fazerem os preparativos.
\VS{53}Mas não o receberam, porque seu rosto demonstrava que ele ia para Jerusalém.
\VS{54}Quando seus discípulos, Tiago e João, viram
{\ADD{isso}} , disseram: “Senhor, queres que digamos que desça fogo do céu e os consuma, como Elias também fez?”
\VS{55}Porém ele se virou, repreendeu-os, e disse: Vós não sabeis de que espírito sois.
\VS{56}Porque o Filho do homem não veio para destruir as almas dos seres humanos, mas sim para salvá-las.E foram para outra aldeia.
\VS{57}E aconteceu que, enquanto eles iam pelo caminho, alguém lhe disse: “Senhor, eu te seguirei para onde quer que fores”.
\VS{58}E Jesus lhe disse: “As raposas têm tocas, e as aves do céu, ninhos; mas o Filho do homem não tem onde deitar a cabeça”.
\VS{59}E disse a outro: “Segue-me”. Porém ele disse: “Senhor, deixa-me que primeiro eu vá, e eu enterre meu pai”.
\VS{60}Mas Jesus lhe disse: “Deixa os mortos enterrarem os seus próprios mortos; tu, porém, vai, e anuncia o reino de Deus”.
\VS{61}E outro também disse: “Senhor, eu te seguirei; mas deixa-me primeiro despedir dos que estão em minha casa”.
\VS{62}E Jesus lhe disse: “Ninguém que colocar a sua mão no arado e olhar para trás é apto para o reino de Deus”.

\par }\Chap{10}{\PP \VerseOne{1}E depois disso, o Senhor ordenou ainda outros setenta, e os mandou de dois em dois adiante de sua face, para toda cidade e lugar aonde ele havia de vir.
\VS{2}E lhes dizia: “A colheita verdadeiramente é grande, mas os trabalhadores são poucos; portanto rogai ao Senhor da colheita para que ele envie trabalhadores para a sua colheita.
\VS{3}Ide; eis que eu vos mando como cordeiros no meio dos lobos.
\VS{4}Não leveis bolsa, nem sacola, nem sandálias; e a ninguém saudeis pelo caminho.
\VS{5}E em qualquer casa que entrardes, dizei primeiro: Paz
{\ADD{seja}} nesta casa.
\VS{6}E se ali houver algum filho da paz, a vossa paz repousará sobre ele; e se não, ela voltará para vós mesmos.
\VS{7}E ficai na mesma casa, comendo e bebendo do que eles vos derem; pois o trabalhador é digno do seu salário. Não vos mudeis de casa em casa.
\VS{8}E em qualquer cidade que entrardes, e vos receberem, comei do que puserem diante de vós.
\VS{9}E curai os enfermos que nela houver, e dizei-lhes: Chegado é para vós o reino de Deus.
\VS{10}Mas em qualquer cidade em que entrardes e não vos receberem, saí pelas ruas, e dizei:
\VS{11}Até o pó da vossa cidade que ficou em nós, sacudimos sobre vós; porém disto sabeis, que o reino de Deus é chegado até vós.
\VS{12}E eu vos digo, que mais tolerável será naquele dia para Sodoma, do que para aquela cidade.
\VS{13}Ai de ti, Corazim! Ai de ti, Betsaida! Porque se em Tiro e em Sidom tivessem sido feitas as maravilhas que foram feitas entre vós, há muito tempo que teriam se arrependido em saco e em cinza.
\VS{14}Portanto, para Tiro e Sidom será mais tolerável no juízo, do que para vós.
\VS{15}E tu, Cafarnaum, que
{\ADD{pensas
}} estar elevada ao céu, até o Xeol
\FTNT{*}{{\FR 10:15 }
Xeol é o lugar dos mortos
} serás derrubada!
\VS{16}Quem vos ouve, ouve a mim; e quem vos rejeita, rejeita a mim; e quem me rejeita, rejeita aquele que me enviou.
\VS{17}E os setenta voltaram com alegria, dizendo: Senhor, até os demônios se sujeitam a nós por teu nome.
\VS{18}E disse-lhes: Eu vi a Satanás, que caía do céu como um raio.
\VS{19}Eis que vos dou poder para pisar sobre serpentes e escorpiões, e sobre toda a força do inimigo, e nada vos fará dano nenhum.
\VS{20}Mas não vos alegreis de que os espíritos se sujeitem a vós; em vez disso, alegrai-vos por vossos nomes estarem escritos nos céus.
\VS{21}Naquela hora Jesus se alegrou em espírito, e disse: Graças te dou, o Pai, Senhor do céu e da terra; porque tu escondeste estas coisas aos sábios e instruídos, e as revelaste às crianças. Sim, Pai, porque assim lhe agradou diante de ti.
\VS{22}Todas as coisas me foram entregues pelo meu Pai; e ninguém sabe quem é o Filho, a não ser o Pai; nem quem é o Pai, a não ser o Filho, e a quem o Filho o quiser revelar.
\VS{23}E virando-se para seus discípulos, disse
{\ADD{-lhes}} à parte: Bem-aventurados os olhos que veem o que vós vedes.
\VS{24}Porque vos digo, que muitos profetas e reis desejaram ver o que vós vedes, e não o viram; o ouvir o que vós ouvis, e não o ouviram.
\VS{25}E eis que um certo estudioso da Lei se levantou, tentando-o, e dizendo: Mestre, o que devo fazer para ter para herdar a vida eterna?
\VS{26}E ele lhe disse: O que está escrito na Lei? Como tu
{\ADD{a}} lês?
\VS{27}E respondendo ele, disse: Amarás ao Senhor teu Deus de todo teu coração, e de toda tua alma, e de todas tuas forças, e de todo teu entendimento; e
{\ADD{amarás}} a teu próximo como a ti mesmo.
\VS{28}E disse-lhe: Respondeste bem; faze isso, e viverás.
\VS{29}Mas ele, querendo se justificar, disse a Jesus: E quem é o meu próximo?
\VS{30}E respondendo Jesus, disse: Um homem descia de Jerusalém a Jericó, e foi atacado por assaltantes, que também tiraram suas roupas, espancaram-no, e se foram, deixando-o meio morto.
\VS{31}E por acaso descia um certo sacerdote pelo mesmo caminho, e vendo-o, passou longe
{\ADD{dele}} .
\VS{32}E semelhantemente também um levita, chegando junto a aquele lugar, veio, e vendo-o, passou longe
{\ADD{dele}}
\VS{33}Porém um certo samaritano, que ia pelo caminho, veio junto a ele, e vendo-o, teve compaixão
{\ADD{dele}} .
\VS{34}E chegando-se, amarrou-lhe
{\ADD{um curativo}} nas feridas, pondo-lhe nelas azeite e vinho; e pondo-o sobre o animal que o transportava, levou-o para uma hospedaria, e cuidou dele.
\VS{35}E partindo-se no outro dia, tirou dois dinheiros, e os deu para o hospedeiro; e disse-lhe: Cuide dele; e tudo o que gastares a mais, eu te pagarei quando voltar.
\VS{36}Quem, pois, destes três te parece que foi o próximo daquele que foi atacado por assaltantes?
\VS{37}Ele disse: Aquele que agiu tendo misericórdia com ele.Então Jesus lhe disse: Vai, e faze da mesma maneira.
\VS{38}E aconteceu que eles, enquanto eles caminhavam, ele entrou em uma aldeia; e uma certa mulher, de nome Marta, o recebeu em sua casa.
\VS{39}E esta tinha uma irmã, chamada Maria, a qual, sentando-se também aos pés de Jesus, ouvia sua palavra.
\VS{40}Marta, porém, ficava muito ocupada com muitos serviços; e ela, vindo, disse: Senhor, não te importas que minha irmã me deixe sozinha para servir? Dize a ela, pois, que me ajude.
\VS{41}E respondendo Jesus, disse-lhe: Marta, Marta,
{\ADD{tu és}} preocupada com muitas coisas, e perturbada por elas;
\VS{42}Mas
{\ADD{somente}} uma coisa é necessária. E Maria escolheu a parte boa, a qual não lhe será tirada.

\par }\Chap{11}{\PP \VerseOne{1}E aconteceu que ele estava orando em um certo lugar. Quando terminou, lhe disse um de seus discípulos: Senhor, ensina-nos a orar, como João também ensinou a seus discípulos.
\VS{2}E ele lhes disse: Quando orardes, dizei: Pai nosso, que
{\ADD{estás}} nos céus, santificado seja o teu nome; venha o teu Reino; seja feita a tua vontade,
{\ADD{assim}} na terra como no céu.
\VS{3}Dá-nos cada dia nosso pão diário.
\VS{4}E perdoa-nos nossos pecados, pois também perdoamos a todo aquele que nos deve. E não nos ponhas em tentação, mas livra-nos do mal.
\VS{5}Disse-lhes também: Qual de vós que, tendo um amigo, se for a ele à meia-noite, e lhe disser: Amigo, empresta-me três pães.
\VS{6}Porque um amigo meu veio de viagem até mim, e nada tenho para lhe apresentar.
\VS{7}E ele de dentro, respondendo, disser: Não me perturbe! A porta já está fechada, e meus filhos estão comigo na cama; não posso me levantar para te dar.
\VS{8}Digo-vos, que ainda que ele não se levante para lhe dar por seu seu amigo, contudo, por sua teimosia ele se levantará, e lhe dará tudo quanto ele precisar.
\VS{9}E eu vos digo: pedi, e será vos dado; buscai, e achareis; batei, e vos será aberto.
\VS{10}Porque todo o que pede, recebe; e quem busca, acha; e quem bate, lhe será aberto.
\VS{11}E que pai, dentre vós, a quem o filho pedir pão, lhe dará uma pedra? Ou, se também
{\ADD{lhe pedir}} peixe, no lugar do peixe lhe dará uma serpente?
\VS{12}Ou se também pedir um ovo, lhe dará um escorpião?
\VS{13}Pois se vós, sendo maus, sabeis dar bons presentes para vossos filhos, quanto mais dará
{\ADD{vosso}} Pai celestial, o Espírito Santo, a aqueles que lhe pedirem?
\VS{14}E
{\ADD{Jesus}} estava expulsando um demônio, e
{\ADD{este}} era mudo. E aconteceu que, saindo o demônio, o mudo falou, e as multidões se maravilharam.
\VS{15}Porém alguns deles diziam: Ele expulsa aos demônios por Belzebu, príncipe dos demônios!
\VS{16}E outros, tentando
{\ADD{-o}} ,pediam-lhe um sinal do céu.
\VS{17}Mas ele, conhecendo seus pensamentos, disse-lhes: Todo reino dividido contra si mesmo é transformado num vazio, e a casa contra casa cai.
\VS{18}E se também Satanás está dividido contra si mesmo, como durará o seu reino? Porque dizeis: Por Belzebu ele expulsa aos demônios.
\VS{19}E se eu expulso aos demônios por Belzebu, por quem os seus filhos expulsam? Portanto eles serão vossos juízes.
\VS{20}Mas se eu expulso aos demônios pelo dedo de Deus, portanto o Reino de Deus chegou a vós.
\VS{21}Quando o valente, armado, guarda seu palácio, seus bens estão em paz.
\VS{22}Mas vindo outro mais valente que ele, e vencendo-o, toma
{\ADD{-lhe}} toda sua armadura, em que confiava, e reparte seus despojos.
\VS{23}Quem não é comigo, é contra mim; e quem comigo não ajunta, espalha.
\VS{24}Quando o espírito imundo tem saído da pessoa, ele anda por lugares secos, buscando repouso; e não
{\ADD{o}} achando, diz: Voltarei para minha casa, de onde saí.
\VS{25}E vindo, acha-a varrida e adornada.
\VS{26}Então vai, e toma consigo outros sete espíritos piores que ele, e entrando, habitam ali; e as últimas coisas de tal pessoa são piores que as primeiras.
\VS{27}E aconteceu que, dizendo ele estas coisas, uma mulher da multidão, levantando a voz, lhe disse: Bem-aventurado o ventre que te trouxe, e os peitos que mamaste!
\VS{28}Mas ele disse: Antes bem-aventurados os que ouvem a palavra de Deus, e a guardam.
\VS{29}E ajuntando as multidões, começou a dizer: Maligna é esta geração; busca sinal, mas sinal não lhe será dado, a não ser o sinal de Jonas o profeta.
\VS{30}Porque como Jonas foi sinal para os ninivitas, assim também será o Filho do homem para esta geração.
\VS{31}A rainha do Sul se levantará em juízo com as pessoas desta geração, e as condenará; pois até dos fins da terra veio para ouvir a sabedoria de Salomão; e eis que mais que Salomão está aqui.
\VS{32}Os homens de Nínive se levantarão em juízo com esta geração, e a condenarão; pois com a pregação de Jonas se converteram; e eis que mais que Jonas está aqui.
\VS{33}E ninguém, acendendo a lâmpada,
{\ADD{a}} põe em
{\ADD{lugar}} oculto, nem debaixo da caixa, mas na luminária, para que os que entrarem vejam a luz.
\VS{34}A lâmpada do corpo é o olho. Sendo pois teu olho bom, também todo teu corpo será luminoso; porém se for mau, também
{\ADD{todo}} teu corpo será tenebroso.
\VS{35}Olha pois que a luz que em ti há não sejam trevas.
\VS{36}Então se sendo teu corpo todo luminoso, não tendo parte alguma escura, ele todo será iluminado, como quando a lâmpada com
{\ADD{seu}} brilho te ilumina.
\VS{37}E estando ele
{\ADD{ainda}} falando, um fariseu lhe rogou que viesse para jantar com ele; e entrando, sentou-se
{\ADD{à mesa}} ;
\VS{38}E vendo
{\ADD{-o}} o fariseu, maravilhou-se de que não tinha se lavado antes de jantar.
\VS{39}E o Senhor lhe disse: Agora vós, os fariseus, limpais o exterior do copo e do prato; porém vosso interior está cheio de roubo e maldade.
\VS{40}Loucos, o que fez o exterior não fez também o interior?
\VS{41}Porém daí de esmola o que tendes; e eis que tudo vos será limpo.
\VS{42}Mais ai de vós, fariseus, que dizimais a hortelã, e a arruda, e toda hortaliça; e pelo juízo e amor de Deus passais longe. Estas coisas era necessário fazer, e não deixar as outras.
\VS{43}Ai de vós, fariseus, que amais os primeiros assentos nas sinagogas, e as saudações nas praças.
\VS{44}Ai de vós, escribas e fariseus, hipócritas, que sois como as sepulturas disfarçadas, e as pessoas que andam sobre elas não sabem.
\VS{45}E respondendo um dos estudiosos da Lei, disse-lhe: Mestre, quando dizes isto também afrontas a nós.
\VS{46}Porém ele disse: Ai de vós também, estudiosos da Lei, que carregais as pessoas com cargas pesadas para levar, e vós mesmos nem ainda com um de vossos dedos tais cargas tocais.
\VS{47}Ai de vós, que construís os sepulcros dos profetas, e vossos pais os mataram.
\VS{48}Bem testemunhais pois, que também consentis nas obras de vossos pais; porque eles os mataram, e vós edificais seus sepulcros.
\VS{49}Portanto também diz a sabedoria de Deus: Profetas e apóstolos lhes mandarei; e deles a
{\ADD{uns}} matarão, e a
{\ADD{outros}} lançarão fora;
\VS{50}Para que desta geração seja requerido o sangue de todos os profetas, que foi derramado desde a fundação do mundo;
\VS{51}Desde o sangue de Abel, até o sangue de Zacarias, que foi morto entre o altar e a casa
{\ADD{de Deus}} ; assim vos digo, será requerido desta geração.
\VS{52}Ai de vós, estudiosos da Lei, que tomastes a chave do conhecimento; vós mesmos não entrastes, e impedistes aos que estavam entrando.
\VS{53}E dizendo-lhes estas coisas, os escribas e os fariseus começaram a apertá-lo fortemente, e tentar lhe fazer falar de muitas coisas,
\VS{54}armando-lhe ciladas, e procurando caçar alguma coisa de sua boca, para o poderem acusar.

\par }\Chap{12}{\PP \VerseOne{1}Juntando-se, entretanto, muitos milhares da multidão, tanto que se atropelavam uns aos outros, começou a dizer primeiramente a seus discípulos: Guardai-vos do fermento dos fariseus, que é hipocrisia.
\VS{2}E nada há encoberto que não haja de ser descoberto; nem oculto que não haja de ser sabido.
\VS{3}Portanto tudo o que dissestes nas trevas, será ouvido na luz; e o que falastes ao ouvido nos quartos, será pregado sobre os telhados.
\VS{4}E digo-vos, amigos meus, não temais aos que matam o corpo, e depois não têm mais o que possam fazer.
\VS{5}Mas eu vos mostrarei a quem deveis temer; temei a aquele, que depois de matar,
{\ADD{também}} tem poder para lançar no inferno; sim, a este temei.
\VS{6}Não se vendem cinco passarinhos por duas pequenas moedas? E nem um deles está esquecido diante de Deus.
\VS{7}E até os cabelos de vossa cabeça estão todos contados; não temais, pois; mais valeis vós que muitos passarinhos.
\VS{8}E vos digo que todo aquele que me confessar diante dos seres humanos, também o Filho do homem o confessará diante dos anjos de Deus.
\VS{9}Mas quem me negar diante dos seres humanos será negado diante dos anjos de Deus.
\VS{10}E a todo aquele que disser
{\ADD{alguma}} palavra contra o Filho do homem, lhe será perdoado, mas ao que blasfemar contra o Espírito Santo, não lhe será perdoado.
\VS{11}E quando vos trouxerem às sinagogas, aos magistrados e autoridades, não estejais ansiosos, como, ou que, em
{\ADD{vossa}} defesa deveis dizer, ou que deveis falar.
\VS{12}Porque na mesma hora o Espírito Santo vos ensinará o que deveis falar.
\VS{13}E um da multidão lhe disse: Mestre, dize a meu irmão que reparta a herança comigo.
\VS{14}Mas ele lhe disse: Homem, quem me pôs por juiz, ou repartidor sobre vós?
\VS{15}E disse-lhes: Olhai, e tomai cuidado com a ganância; porque a vida de alguém não consiste na abundância dos bens que possui.
\VS{16}E propôs-lhes uma parábola, dizendo: A terra de um homem rico tinha frutificado bem.
\VS{17}E ele questionava a si mesmo, dizendo: Que farei? Porque não tenho onde juntar meus frutos.
\VS{18}E disse: Farei isto: derrubarei meus celeiros, e construirei maiores, e ali juntarei toda esta minha colheita, e estes meus bens.
\VS{19}E direi à minha alma: Alma, muitos bens tens guardados, para muitos anos; descansa, come, bebe, alegra-te!
\VS{20}Porém Deus lhe disse: Louco, esta noite te pedirão tua alma; e o que tens preparado, de quem será?
\VS{21}Assim
{\ADD{é}} o que junta tesouros para si, mas não é rico em Deus.
\VS{22}E disse a seus discípulos: Portanto vos digo, não estejais ansiosos por vossa vida, que comereis; nem pelo corpo, que vestireis.
\VS{23}Mais é a vida que o alimento, e
{\ADD{mais}} o corpo que o vestido.
\VS{24}Considerai os corvos, que nem semeiam, nem ceifam; nem tem armazém, nem celeiro; e Deus os alimenta.
\VS{25}E quem de vós pode, com sua ansiedade, acrescentar um côvado à sua altura?
\VS{26}Pois, se não podeis nem mesmo com algo pequeno, por que estais ansiosos com o resto?
\VS{27}Considerai os lírios, como crescem; não trabalham, nem fiam; e digo-vos, que nem mesmo Salomão, em toda sua glória,
{\ADD{chegou}} a se vestir como um deles.
\VS{28}E se assim Deus veste a erva, que hoje está no campo, e amanhá é lançada no forno, quanto mais
{\ADD{vestirá}} a vós, homens de pouca fé?
\VS{29}Vós, pois, não pergunteis que comereis, ou que bebereis; e não andeis preocupados.
\VS{30}Porque todas estas coisas, os gentios do mundo as buscam; mas vosso Pai sabe que necessitais destas coisas.
\VS{31}Mas buscai o Reino de Deus, e todas estas coisas vos serão acrescentadas.
\VS{32}Não temas, ó pequeno rebanho; porque vosso Pai se agradou de dar a vós o Reino.
\VS{33}Vendei o que tendes, e daí esmola. Fazei para vós bolsas que não se envelheçam; tesouro nos céus que nunca se deprecia; onde ladrão não chega, nem a traça destrói.
\VS{34}Porque aonde estiver vosso tesouro, ali estará também vosso coração.
\VS{35}Estejam devidamente vestidos vossos quadris, e acesas as lâmpadas.
\VS{36}E sede vós semelhantes às pessoas que esperam a seu senhor quando voltar do casamento; para que quando ele vier, e bater, logo possam lhe abrir.
\VS{37}Bem-aventurados aqueles servos, os quais, quando o Senhor vier,
{\ADD{os}} achar vigiando; em verdade vos digo que ele se vestirá, e os fará se sentarem
{\ADD{à mesa}} ,e chegando-se, os servirá.
\VS{38}E ainda que venha à segunda vigília; e
{\ADD{que}} venha a terceira vigília, e assim
{\ADD{os}} achar, bem-aventurados são tais servos.
\VS{39}Isto, porém, sabei: que se o chefe da casa soubesse à que hora o ladrão viria, ele vigiaria, e não deixaria sua casa sofrer dano.
\VS{40}Vós, pois, também estejais prontos; porque o Filho do homem virá à hora que não imaginais.
\VS{41}E Pedro lhe disse: Senhor, dizes tu esta parábola para nós, ou também para todos?
\VS{42}E o Senhor disse: Qual é, pois, o mordomo fiel e prudente, a quem
{\ADD{seu}} senhor puser sobre seus servos, para que
{\ADD{lhes}} dê alimento no tempo certo?
\VS{43}Bem-aventurado aquele servo ao qual, quando seu senhor vier, o achar fazendo assim.
\VS{44}Em verdade vos digo, que o porá sobre todos os seus bens.
\VS{45}Mas se aquele servo disser em seu coração: Meu senhor está demorando para vir; e começar a espancar aos servos e servas, e a comer, e a beber, e a se embebedar,
\VS{46}Virá o senhor daquele servo, no dia em que ele não espera, e na hora que ele não sabe; e será partido em dois, e porá sua porção com os incrédulos.
\VS{47}E o servo que sabia a vontade de seu senhor, e não se preparou, nem fez conforme a sua vontade, será muito espancado.
\VS{48}Mas o que não sabia, e fez
{\ADD{coisas}} dignas de pancadas, será pouco espancado. E a qualquer que muito for dado, muito se lhe pedirá, e ao que muito se confiou, muito mais lhe será exigido.
\VS{49}Eu vim para lançar fogo na terra; e o que mais posso querer, se já está aceso?
\VS{50}Porém há um batismo que tenho que ser batizado; e como me angustio até que se venha a cumprir!
\VS{51}Vós pensais que vim para dar paz à terra? Não, eu vos digo; mas antes
{\ADD{vim para trazer}} divisão.
\VS{52}Porque daqui em diante cinco estarão divididos em uma casa; três contra dois, e dois contra três.
\VS{53}O pai estará dividido contra o filho, e o filho contra o pai; a mãe contra a filha, e a filha contra a mãe; a sogra contra sua nora, e a nora contra sua sogra.
\VS{54}E ele dizia também para as multidões: Quando vedes a nuvem que vem do ocidente, logo dizeis: Lá vem chuva; E assim acontece.
\VS{55}E quando venta do sul, dizeis: Haverá calor; E assim acontece.
\VS{56}Hipócritas! Sabeis entender a aparência da terra e do céu; e como não entendeis este tempo?
\VS{57}E por que também não julgais por vós mesmos o que é justo?
\VS{58}Pois quando fores com teu adversário à autoridade, procura te livrares dele no caminho, para que ele não venha a te levar ao juiz, e o juiz te entregue ao oficial de justiça, e o oficial de justiça te lance na prisão.
\VS{59}Eu te digo que não sairás dali enquanto não pagares até a última moeda.
\FTNT{*}{{\FR 12:59 }
moeda
lit. quadrante, o nome de uma moeda de pequeno valor
}

\par }\Chap{13}{\PP \VerseOne{1}E naquele mesmo tempo estavam ali presentes alguns, que lhe contavam dos galileus cujo sangue Pilatos tinha misturado com seus sacrifícios.
\VS{2}E respondendo Jesus, disse-lhes: Vós pensais que estes galileus foram mais pecadores, por terem sofrido estas coisas?
\VS{3}Não, eu vos digo; antes, se vós não vos arrependerdes, todos de modo semelhante perecereis.
\VS{4}Ou aqueles dezoito, sobre os quais a torre em Siloé caiu, matando-os; pensais que eram mais culpados dos que todos as pessoas que moravam em Jerusalém?
\VS{5}Não, eu vos digo; antes, se vós não vos arrependerdes, todos de modo semelhante perecereis.
\VS{6}E dizia esta parábola: Um certo
{\ADD{homem}} tinha uma figueira plantada em sua vinha, e veio até ela para buscar fruto, e não achou.
\VS{7}E disse ao que cuidava da vinha: Eis que há três anos, que venho para buscar fruto nesta figueira, e não
{\ADD{o}} acho; corta-a; por que ainda ocupa inutilmente a terra?
\VS{8}E respondendo ele, disse-lhe: Senhor, deixa-a
{\ADD{ainda}} este ano, até que eu a escave ao redor, e a adube;
\VS{9}E se der fruto,
{\ADD{deixa-a ficar}} ; se não, tu a cortarás depois.
\VS{10}E ensinava em uma das sinagogas num sábado.
\VS{11}E eis que estava ali uma mulher, que havia dezoito anos que tinha um espírito de enfermidade; e andava encurvada, e de maneira nenhuma ela podia se endireitar.
\VS{12}E Jesus vendo-a, chamou-a para si, e disse-lhe: Mulher, livre estás de tua enfermidade.
\VS{13}E pôs as mãos sobre ela, e logo ela se endireitou, e glorificava a Deus.
\VS{14}E o chefe da sinagoga, irritado por Jesus ter curado no sábado, respondendo, disse à multidão: Há seis dias em que se deve trabalhar; nestes dias, pois, vinde para ser curados, e não no dia de sábado.
\VS{15}Porém o Senhor lhe respondeu, e disse: Hipócrita, no sábado cada um de vós não desata seu boi ou jumento da manjedoura, e o leva para dar de beber?
\VS{16}E não convinha soltar desta ligadura no dia de sábado a esta
{\ADD{mulher}} ,que é filha de Abraão, a qual, eis que Satanás já a havia ligado há dezoito anos?
\VS{17}E ele, dizendo estas coisas, todos seus adversários ficaram envergonhados; e todo o povo se alegrava de todas as coisas gloriosas que eram feitas por ele.
\VS{18}E dizia: A que o Reino de Deus é semelhante? E a que eu o compararei?
\VS{19}Semelhante é ao grão da mostarda, que um homem, tomando-o, lançou-o em sua horta; e cresceu, e fez-se uma grande árvore, e as aves dos céus fizeram ninhos em seus ramos.
\VS{20}E disse outra vez: A que compararei o Reino de Deus?
\VS{21}Semelhante é ao fermento, que a mulher, tomando-o, o escondeu em três medidas de farinha, até tudo ficar levedado.
\VS{22}E andava de cidade em cidade, e de aldeia em aldeia ensinando, e caminhando para Jerusalém.
\VS{23}E disse-lhe um: Senhor, são também poucos os que se salvam? E ele lhes disse:
\VS{24}Trabalhai para entrar pela porta estreita; porque eu vos digo,
{\ADD{que}} muitos procuraram entrar, e não puderam.
\VS{25}Porque quando o chefe da casa se levantar, e fechar a porta, e se estiverdes de fora, e começardes a bater à porta, dizendo: Senhor, senhor, abre-nos! E respondendo ele, vos disser: Não vos conheço,
{\ADD{nem sei}} de onde vós sois.
\VS{26}Então começareis a dizer: Em tua presença temos comido e bebido, e tens ensinado em nossas ruas.
\VS{27}E ele dirá: Digo-vos que não vos conheço,
{\ADD{nem sei}} de onde vós sois; afastai-vos de mim, vós todos praticantes de injustiça.
\VS{28}Ali haverá choro e ranger de dentes, quando virdes a Abraão, a Isaque, a Jacó, e a todos os profetas no Reino de Deus; mas vós
{\ADD{sendo}} lançados fora.
\VS{29}E virão
{\ADD{pessoas}} do oriente, e do ocidente, e do norte, e do sul, e se sentarão
{\ADD{à mesa}} no Reino de Deus.
\VS{30}E eis que há
{\ADD{alguns}} dos últimos que serão primeiros, e há
{\ADD{alguns}} dos primeiros que serão últimos.
\VS{31}Naquele mesmo dia, chegaram uns fariseus, dizendo-lhe: Sai, e vai-te daqui, porque Herodes quer te matar.
\VS{32}E disse-lhes: Ide, e dizei a aquela raposa: eis que expulso demônios, e faço curas hoje e amanhã, e ao terceiro dia eu terei completado.
\VS{33}Porém é necessário que hoje, e amanhã, e no
{\ADD{dia}} seguinte eu caminhe; porque um profeta não pode morrer fora de Jerusalém.
\VS{34}Jerusalém, Jerusalém, que matas aos profetas, e apedrejas aos que te são enviados: quantas vezes eu quis juntar teus filhos, como a galinha junta seus pintos debaixo de suas asas, e não quisestes?
\VS{35}Eis que vossa casa é deixada deserta para vós. E em verdade vos digo, que não me vereis até que venha
{\ADD{o tempo}} em que digais: Bendito aquele que vem no nome do Senhor.

\par }\Chap{14}{\PP \VerseOne{1}E aconteceu que, entrando ele num sábado para comer pão na casa de um dos chefes dos fariseus, eles estavam o observando.
\VS{2}E eis que um certo homem com o corpo inchado estava ali diante dele.
\VS{3}E respondendo Jesus, falou aos estudiosos da Lei, e aos fariseus, dizendo: É lícito curar no sábado?
\VS{4}Porém eles ficaram calados; e ele, tomando
{\ADD{-o}} ,o curou, e
{\ADD{o}} despediu.
\VS{5}E ele, respondendo-lhes, disse: De qual que vós cairá o jumento, ou o boi em algum poço que,
{\ADD{mesmo}} no sábado, não o tire logo?
\VS{6}E nada podiam lhe responder a estas coisas.
\VS{7}E vendo como os convidados escolhiam os primeiros assentos, disse-lhes uma parábola:
\VS{8}Quando fores convidados para o casamento de alguém, não te sentes no primeiro assento, para que não
{\ADD{aconteça de}} se outro convidado mais digno que tu estiver,
\VS{9}E venha o que convidou a ti e a ele, e te diga: Dá lugar a este; E então, com vergonha, tenhas que tomar o último lugar.
\VS{10}Mas quando fores convidado, vai, e senta-te no último lugar; para que quando vier o que te convidou, te diga: Amigo, sobe para
{\ADD{este assento}} melhor. Então terás honra diante dos que estiverem sentados contigo
{\ADD{à mesa}} .
\VS{11}Porque qualquer que exaltar a si mesmo, e aquele que humilhar a si mesmo, será exaltado.
\VS{12}E dizia também ao que tinha lhe convidado: Quando fizeres um jantar, ou uma ceia, não chames a teus amigos, nem a teus irmãos, nem a teus parentes, nem a
{\ADD{teus}} vizinhos ricos, para que eles também em algum tempo não te convidem de volta, e tu sejas recompensado.
\VS{13}Mas quando fizeres convite, chama aos pobres, aleijados, mancos
{\ADD{e}} cegos.
\VS{14}E serás bem-aventurado, porque não eles não têm como te recompensar; pois tu serás recompensado na ressurreição dos justos.
\VS{15}E um dos que juntamente estavam sentados
{\ADD{à mesa}} , ouvindo isto, disse-lhe: Bem-aventurado aquele que comer pão no Reino de Deus.
\VS{16}Porém ele lhe disse: Um certo homem fez um grande jantar, e convidou a muitos.
\VS{17}E na hora do jantar, mandou seu servo para dizer aos convidados: Vinde, que tudo já está preparado.
\VS{18}E cada um deles todos começou a dar desculpas. O primeiro lhe disse: Comprei um campo, e tenho que ir vê-lo; peço-te desculpas.
\VS{19}E outro disse: Comprei cinco pares de bois, e vou testá-los; peço-te desculpas.
\VS{20}E outro disse: Casei-me
{\ADD{com}} uma mulher, e portanto não posso vir.
\VS{21}E aquele servo, ao voltar, anunciou estas coisas a seu senhor. Então o chefe da casa, irritado, disse a seu servo: Sai depressa pelas ruas e praças da cidade, e traze aqui aos pobres, e aleijados, e mancos e cegos.
\VS{22}E o servo disse: Senhor, está feito como mandaste, e ainda há lugar.
\VS{23}E o senhor disse ao servo: Sai pelos caminhos, e trilhas, e força-os a entrar, para que minha casa se encha.
\VS{24}Porque eu vos digo, que nenhum daqueles homens que foram convidados experimentará da minha ceia.
\VS{25}E muitas multidões iam com ele; e virando-se, disse-lhes:
\VS{26}Se alguém vier a mim, e não odiar a seu pai, e mãe, e mulher, e filhos, e irmãos, e irmãs, e ainda também sua própria vida, não pode ser meu discípulo.
\VS{27}E qualquer que não levar sua cruz, e vier após mim, não pode ser meu discípulo.
\VS{28}Porque qual de vós, querendo edificar uma torre, não se senta primeiro para fazer as contas dos gastos,
{\ADD{para ver}} se tem o
{\ADD{suficiente}} para a completar?
\VS{29}Para que não aconteça que, depois de ter posto seu fundamento, e não podendo a completar, comecem a escarnecer dele todos os que
{\ADD{o}} virem,
\VS{30}Dizendo: Este homem começou a construir, e não pôde terminar.
\VS{31}Ou qual rei, indo a guerra para lutar contra outro rei, não se senta primeiro para consultar, se pode ir ao encontro com dez mil
{\ADD{soldados}} , vindo contra ele vinte mil?
\VS{32}Se não
{\ADD{puder}} , estando o outro ainda longe, manda
{\ADD{-lhe}} representantes diplomáticos, e roga pela paz.
\VS{33}Assim, portanto, qualquer de vós que não renuncia a tudo, não pode ser meu discípulo.
\VS{34}Bom é o sal; porém se o sal perder o sabor, com o que ele será temperado?
\VS{35}Nem para a terra, nem para adubo serve; lançam-no fora. Quem tem ouvidos para ouvir, ouça.

\par }\Chap{15}{\PP \VerseOne{1}E chegavam-se a ele todos os publicanos e pecadores para o ouvirem.
\VS{2}E os fariseus e escribas murmuravam, dizendo: Este recebe aos pecadores, e come com eles.
\VS{3}E ele lhes propôs esta parábola, dizendo:
\VS{4}Quem de vós, tendo cem ovelhas, e perdendo uma delas, não deixa no deserto as noventa e nove, e vai em
{\ADD{busca}} da perdida, até que a encontre?
\VS{5}E encontrando-a,
{\ADD{não}} a ponha sobre seus ombros, com alegria?
\VS{6}E vindo para casa,
{\ADD{não}} convoque aos amigos, e vizinhos, dizendo-lhes: Alegrai-vos comigo, porque já achei minha ovelha perdida?
\VS{7}Digo-vos, que assim haverá mais alegria no céu por um pecador que se arrepende, do que por noventa e nove justos que não precisam de arrependimento.
\VS{8}Ou que mulher, tendo dez moedas de prata, se perder a uma moeda, não acende a lâmpada, e varre a casa, e busca cuidadosamente até
{\ADD{a}} achar?
\VS{9}E achando
{\ADD{-a}} ,
{\ADD{não}} chame as amigas e as vizinhas, dizendo: Alegrai-vos comigo, porque já achei a moeda perdida!
\VS{10}Assim vos digo, que há alegria diante dos anjos de Deus por um pecador que se arrepende.
\VS{11}E disse: Um certo homem tinha dois filhos.
\VS{12}E disse o mais jovem deles ao pai: Pai, dá-me a parte dos bens que
{\ADD{me}} pertencem. E ele lhe repartiu os bens.
\VS{13}E depois de não muitos dias, o filho mais jovem, juntando tudo, partiu-se para uma terra distante, e ali desperdiçou seus bens, vivendo de forma irresponsável.
\VS{14}E ele, tendo já gastado tudo, houve uma grande fome naquela terra, e ele começou a sofrer necessidade.
\VS{15}E foi, se chegou a um dos cidadãos daquela terra; e
{\ADD{este}} o mandou a seus campos para alimentar porcos.
\VS{16}E ele ficava com vontade de encher seu estômago com os grãos que os porcos comiam, mas ninguém
{\ADD{as}} dava para ele.
\VS{17}E ele, pensando consigo mesmo, disse: Quantos empregados de meu pai tem pão em abundância, e eu
{\ADD{aqui}} morro de fome!
\VS{18}Eu levantarei, e irei a meu pai, e lhe direi: Pai, pequei contra o céu, e diante de ti.
\VS{19}E já não sou digno de ser chamado teu filho; faze-me como a um de teus empregados.
\VS{20}E levantando-se, foi a seu pai. E quando ainda estava longe, o seu pai o viu, e teve compaixão dele; e correndo, caiu ao seu pescoço, e o beijou.
\VS{21}E o filho lhe disse: Pai, pequei contra o céu, e diante de ti; e já não sou digno de ser chamado teu filho.
\VS{22}Mas o pai disse a seus servos: Trazei a melhor roupa, e o vesti; e ponde um anel em sua mão, e sandálias nos
{\ADD{seus}} pés.
\VS{23}E trazei o bezerro engordado, e o matai; e comamos, e nos alegremos.
\VS{24}Porque este meu filho estava morto, e reviveu; tinha se perdido, e foi achado. E começaram a se alegrar.
\VS{25}E seu filho mais velho estava no campo; e quando veio, chegou perto da casa, ouviu a música, e as danças.
\VS{26}E chamando para si um dos servos, perguntou-lhe: O que era aquilo?
\VS{27}E ele lhe disse: Teu irmão chegou; e teu pai matou o bezerro engordado, porque ele voltou são.
\VS{28}Porém ele se irritou, e não queria entrar. Então o seu pai, saindo, rogava-lhe
{\ADD{que entrasse}} .
\VS{29}Mas
{\ADD{o filho}} respondendo, disse ao pai: Eis que eu te sirvo há tantos anos, e nunca desobedeci tua ordem, e nunca me deste um cabrito, para que eu me alegrasse com meus amigos.
\VS{30}Porém, vindo este teu filho, que gastou teus bens com prostitutas, tu lhe mataste o bezerro engordado.
\VS{31}E ele lhe disse: Filho, tu sempre estás comigo, e todas as minhas coisas são tuas.
\VS{32}Mas era necessário se alegrar e animar; porque este teu irmão estava morto, e reviveu; e tinha se perdido, e foi encontrado.

\par }\Chap{16}{\PP \VerseOne{1}E dizia também a seus discípulos: Havia um certo homem rico, o qual tinha um mordomo; e este lhe foi acusado de fazer perder seus bens.
\VS{2}E ele, chamando-o, disse-lhe: Como ouço isto sobre ti? Presta contas de teu trabalho, porque não poderás mais ser
{\ADD{meu}} mordomo.
\VS{3}E disse o mordomo para si mesmo: O que farei, agora que meu senhor está me tirando o trabalho de mordomo? Cavar eu não posso; mendigar eu tenho vergonha.
\VS{4}Eu sei o que farei, para que, quando eu for expulso do meu trabalho de mordomo, me recebam em suas casas.
\VS{5}E chamando a si a cada um dos devedores de seu senhor, disse ao primeiro: Quanto deves a meu senhor?
\VS{6}E ele disse: Cem medidas de azeite. E disse-lhe: Pega a tua conta, senta, e escreve logo cinquenta.
\VS{7}Depois disse a outro: E tu, quanto deves? E ele disse: Cem volumes de trigo. E disse-lhe: Toma tua conta, e escreve oitenta.
\VS{8}E aquele senhor elogiou o injusto mordomo, por ter feito prudentemente; porque os filhos deste mundo são mais prudentes do que os filhos da luz com esta geração.
\VS{9}E eu vos digo: fazei amigos para vós com as riquezas da injustiça, para que quando vos faltar, vos recebam nos tabernáculos eternos.
\VS{10}Quem é fiel no mínimo, também é fiel no muito; e quem é injusto no mínimo, também é injusto no muito.
\VS{11}Pois se nas riquezas da injustiça não fostes fiéis, quem vos confiará as verdadeiras
{\ADD{riquezas}} ?
\VS{12}E se nas
{\ADD{coisas}} dos outros não fostes fiéis, quem vos dará o que é vosso?
\VS{13}Nenhum servo pode servir a dois senhores; porque ou irá odiar a um, e a amar ao outro; ou irá se achegar a um, e desprezar ao outro. Não podeis servir a Deus e às riquezas.
\VS{14}E os fariseus também ouviram todas estas coisas, eles que eram avarentos. E zombaram dele.
\VS{15}E disse-lhes: Vós sois os que justificais a vós mesmos diante dos seres humanos; mas Deus conhece vossos corações. Porque o que é excelente para os seres humanos é odiável diante de Deus.
\VS{16}A Lei e os profetas
{\ADD{foram}} até João; desde então, o Reino de Deus é anunciado, todo homem
{\ADD{tenta entrar}} nele pela força.
\VS{17}E é mais fácil passar o céu e a terra do que cair um traço de alguma letra da Lei.
\VS{18}Qualquer que deixa sua mulher, e casa com outra, adultera; e qualquer que se casa com a deixada pelo marido,
{\ADD{também}} adultera.
\VS{19}Havia porém um certo homem rico, e vestia-se de púrpura, e de linho finíssimo, e festejava todo dia com luxo.
\VS{20}Havia também um certo mendigo, de nome Lázaro, o qual ficava deitado à sua porta cheio de feridas.
\VS{21}E desejava se satisfazer com as migalhas que caíam da mesa do rico; porém vinham também os cães, e lambiam suas feridas.
\VS{22}E aconteceu que o mendigo morreu, e foi levado pelos anjos para o colo de Abraão. E o rico também morreu, e foi sepultado.
\VS{23}E estando no Xeol
\FTNT{*}{{\FR 16:23 }
Xeol é o lugar dos mortos
} em tormentos, ele levantou seus olhos, e viu a Abraão de longe, e a Lázaro junto dele.
\VS{24}E ele, chamando, disse: Pai Abraão, tem misericórdia de mim, e manda a Lázaro que molhe a ponta de seu dedo na água, e refresque a minha língua; porque estou sofrendo neste fogo.
\VS{25}Porém Abraão disse: Filho, lembra-te que em tua vida recebeste teus bens, e Lázaro do mesmo jeito
{\ADD{recebeu}} males. E agora este é consolado, e tu
{\ADD{és}} atormentado.
\VS{26}E, além de tudo isto, um grande abismo está posto entre nós e vós, para os que quisessem passar daqui para vós não possam; nem também os daí passarem para cá.
\VS{27}E disse ele: Rogo-te, pois, ó pai, que o mandes à casa de meu pai.
\VS{28}Porque tenho cinco irmãos, para que lhes dê testemunho; para que também não venham para este lugar de tormento.
\VS{29}Disse-lhe Abraão: Eles têm a Moisés e aos profetas, ouçam-lhes.
\VS{30}E ele disse: Não, pai Abraão; mas se alguém dos mortos fosse até eles, eles se arrependeriam.
\VS{31}Porém
{\ADD{Abraão}} lhe disse: Se não ouvem a Moisés e aos profetas, também não se deixariam convencer, ainda que alguém ressuscite dos mortos.

\par }\Chap{17}{\PP \VerseOne{1}E disse aos discípulos: É impossível que não venham tentações para o pecado,
\FTNT{*}{{\FR 17:1 }
tentações para o pecado
Lit. meios de tropeço
} mas ai
{\ADD{daquele}} por quem
{\ADD{essas tentações}} vierem!
\VS{2}Melhor lhe seria que lhe atasse ao pescoço uma grande pedra de moinho, e fosse lançado no mar, do que conduzir ao pecado um destes pequenos.
\VS{3}Olhai por vós. E se teu irmão pecar contra ti, repreende-o; e se ele se arrepender, perdoa-lhe.
\VS{4}E se pecar contra ti sete vezes ao dia, e se sete vezes ao dia voltar a ti, dizendo: Estou arrependido. Perdoa-lhe.
\VS{5}E os apóstolos disseram ao Senhor: Acrescenta-nos fé.
\VS{6}E o Senhor disse: Se tivésseis fé como um grão de mostarda, diríeis a esta árvore de amoras: Arranca-te daqui pelas tuas raízes, e planta-te no mar, ela vos obedeceria.
\VS{7}E qual de vós terá um servo, lavrando ou apascentando
{\ADD{gado}} que, voltando do campo, logo lhe diga: Chega, e senta
{\ADD{à mesa}} .
\VS{8}E não lhe diga antes: Prepara-me o jantar, e apronta-te, e serve-me, até que eu tenha comido e bebido; e depois, que tu comas e bebas.
\VS{9}Por acaso
{\ADD{o senhor}} agradece a tal servo, porque fez o que lhe foi mandado? Acho que não.
\VS{10}Assim também vós, quando fizerdes tudo o que vos for mandado, dizei: Somos servos inúteis, porque fizemos
{\ADD{somente}} o que devíamos fazer.
\VS{11}E aconteceu que, indo ele para Jerusalém, passou por meio da Samaria e da Galileia.
\VS{12}E entrando em uma certa aldeia, saíram-lhe ao encontro dez homens leprosos, os quais pararam de longe.
\VS{13}E levantaram a voz, dizendo: Jesus, Mestre, tem misericórdia de nós!
\VS{14}E ele, vendo-os, disse-lhes: Ide, e mostrai-vos aos sacerdotes.E aconteceu que, enquanto eles iam, ficaram limpos.
\VS{15}E vendo um deles que estava são, voltou, glorificando a Deus a alta voz.
\VS{16}E caiu com o rosto a seus pés, agradecendo-lhe; e este era samaritano.
\VS{17}E respondendo Jesus, disse: Não foram os dez limpos? E onde estão os nove?
\VS{18}Não houve quem voltasse para dar glória a Deus, a não ser este estrangeiro?
\VS{19}E disse-lhe: Levanta-te, e vai; tua fé te salvou.
\VS{20}E perguntado pelos fariseus
{\ADD{sobre}} quando o Reino de Deus viria, respondeu-lhes, e disse: O Reino de Deus não vem com aparência visível.
\VS{21}Nem dirão: Eis aqui, ou Eis ali, porque eis que o Reino de Deus está entre vós.
\VS{22}E disse aos discípulos: Dias virão, quando desejareis ver um dos dias do Filho do homem, e não
{\ADD{o}} vereis.
\VS{23}E vos dirão: Eis que
{\ADD{ele está}} aqui, ou Eis que
{\ADD{ele está}} ali, não vades, nem sigais.
\VS{24}Porque como o relâmpago, que relampeja desde o começo do céu, e brilha até ao fim do céu, assim será também o Filho do homem em seu dia.
\VS{25}Mas é necessário primeiro sofrer muito, e ser rejeitado por esta geração.
\VS{26}E como aconteceu nos dias de Noé, assim será também nos dias do Filho do homem.
\VS{27}Comiam, bebiam, se casavam, e se davam em casamento, até o dia em que Noé entrou na arca; e veio o dilúvio, e destruiu a todos.
\VS{28}Como também da mesma maneira aconteceu nos dias de Ló, comiam, bebiam, compravam, vendiam, plantavam,
{\ADD{e}} construíam.
\VS{29}Mas o dia em que Ló saiu de Sodoma, choveu fogo e enxofre do céu, e destruiu a todos.
\VS{30}Assim será
{\ADD{também}} no dia em que o Filho do homem se manifestar.
\VS{31}Naquele dia, o que estiver no telhado, e suas ferramentas em casa, não desça para pegá-las; e o que
{\ADD{estiver}} no campo, não volte para trás.
\VS{32}Lembrai-vos da mulher de Ló.
\VS{33}Qualquer que procurar salvar sua vida a perderá; e qualquer que a perder, irá salvá-la.
\VS{34}Digo-vos que naquela noite, dois estarão em uma cama; um será tomado, e o outro será deixado.
\VS{35}Duas estarão juntas moendo; uma será tomada, e a outra será deixada.
\VS{36}Dois estarão no campo; um será tomado, e o outro será deixado.
\VS{37}E respondendo, disseram-lhe: Onde, Senhor?E ele lhes disse: Onde estiver o corpo, ali os abutres se juntarão.

\par }\Chap{18}{\PP \VerseOne{1}E
{\ADD{Jesus}} lhes disse também uma parábola
{\ADD{sobre}} o dever de sempre orar, e nunca se cansar.
\VS{2}Dizendo: Havia um certo juiz em uma cidade, que não temia a Deus, nem respeitava pessoa alguma.
\VS{3}Havia também naquela mesma cidade uma certa viúva, e vinha até ele, dizendo: Faze-me justiça com meu adversário.
\VS{4}E por um
{\ADD{certo}} tempo ele não quis; mas depois disto, disse para si: Ainda que eu não tema a Deus, nem respeite pessoa alguma,
\VS{5}Porém, porque esta viúva me incomoda, eu lhe farei justiça, para que ela pare de vir me chatear.
\VS{6}E disse o Senhor: Ouvi o que diz o juiz injusto.
\VS{7}E Deus não fará justiça para seus escolhidos, que clamam a ele de dia e de noite? Demorará com eles?
\VS{8}Digo-vos que depressa lhes fará justiça. Porém, quando o Filho do homem vier, por acaso ele achará fé na terra?
\VS{9}E disse também a uns, que tinham confiança de si mesmos que eram justos, e desprezavam aos outros, esta parábola:
\VS{10}Dois homens subiram ao Templo para orar, um fariseu, e o outro publicano.
\VS{11}O fariseu, estando de pé, orava consigo desta maneira: Ó Deus, eu te agradeço, porque não sou como os outros homens, ladrões, injustos e adúlteros; nem
{\ADD{sou}} como este publicano.
\VS{12}Jejuo duas vezes por semana,
{\ADD{e}} dou dízimo de tudo quanto possuo.
\VS{13}E o publicano, estando em pé de longe, nem mesmo queria levantar os olhos ao céu, mas batia em seu peito, dizendo: Ó Deus, tem misericórdia de mim,
{\ADD{que sou}} pecador.
\VS{14}Digo-vos que este desceu mais justificado à sua casa do que aquele outro; porque qualquer que a si mesmo se exalta, será humilhado; e qualquer que a si mesmo se humilha, será exaltado.
\VS{15}E traziam-lhe também crianças pequenas, para que as tocasse; e os discípulos, vendo, os repreendiam.
\VS{16}Mas Jesus, chamando-lhes para si, disse: Deixai as crianças virem a mim, e não as impeçais; porque das tais é o Reino de Deus.
\VS{17}Em verdade vos digo, que qualquer que não receber o Reino de Deus como criança, não entrará nele.
\VS{18}E um certo líder lhe perguntou, dizendo: Bom mestre, o que tenho que fazer para herdar a vida eterna?
\VS{19}E Jesus lhe disse: Por que me chamas de bom? Ninguém
{\ADD{é}} bom, a não ser um: Deus.
\VS{20}Tu sabes os mandamentos: não adulterarás, não matarás, não furtarás, não darás falso testemunho, honra a teu pai e a tua mãe.
\VS{21}E ele disse: Todas estas coisas tenho guardado desde minha juventude.
\VS{22}Porém Jesus, ouvindo isto, disse-lhe: Ainda uma coisa te falta: vende tudo quanto tens, e reparte-o entre os pobres, e terás um tesouro no céu; e vem, segue-me.
\VS{23}Mas ele, ouvindo isto, ficou muito triste, porque era muito rico.
\VS{24}E vendo Jesus que ele tinha ficado muito triste, disse: Como é difícil os que têm
{\ADD{muitos}} bens entrarem no Reino de Deus!
\VS{25}Porque é mais fácil um camelo entrar pelo olho de uma agulha do que um rico entrar no Reino de Deus.
\VS{26}E os que ouviram
{\ADD{isto}} ,disseram: Quem então pode se salvar?
\VS{27}E ele disse: As coisas que são impossíveis para os seres humanos são possíveis para Deus.
\VS{28}E Pedro disse: Eis que deixamos tudo, e temos te seguido.
\VS{29}E ele lhes disse: Em verdade vos digo, que há ninguém que, tenha deixado casa, ou pais, ou irmãos, ou mulher, ou filhos,
\VS{30}Que não venha a receber de volta muito mais nestes tempos, e nos tempos vindouros
{\ADD{receba}} a vida eterna.
\VS{31}E tomando consigo aos doze, disse-lhes: Eis que estamos subindo a Jerusalém, e se cumprirá no Filho do homem tudo o que
{\ADD{está}} escrito pelos profetas.
\VS{32}Porque ele será entregue aos gentios, e escarnecido, injuriado e cuspido.
\VS{33}E
{\ADD{depois de}} açoitá-lo, o matarão; e ao terceiro dia ressuscitará.
\VS{34}E eles nada entendiam destas coisas; e esta palavra lhes era oculta; e não entendiam o que estava sendo lhes dito.
\VS{35}E aconteceu que ele, chegando perto de Jericó, estava um cego sentado junto ao caminho, mendigando.
\VS{36}E este, ouvindo a multidão passar, perguntou: O que era aquilo?
\VS{37}E disseram-lhe que Jesus Nazareno estava passando.
\VS{38}Então clamou, dizendo: Jesus, Filho de Davi, tem misericórdia de mim!
\VS{39}E os que estavam mais a frente o repreendiam, para que calasse; porém ele clamava ainda mais: Filho de Davi, tem misericórdia de mim!
\VS{40}Então Jesus, parando, mandou que o trouxessem para si; e chegando ele, perguntou-lhe,
\VS{41}Dizendo: Que queres que eu te faça?E ele disse: Senhor, que eu veja.
\VS{42}E Jesus lhe disse: Vê, tua fé te salvou.
\VS{43}E logo ele viu, e o seguia, glorificando a Deus. E o todo o povo, vendo
{\ADD{isto}} , dava louvores a Deus.

\par }\Chap{19}{\PP \VerseOne{1}E
{\ADD{Jesus}} entrou e foi passando por Jericó.
\VS{2}E eis que havia ali um homem, chamado pelo nome de Zaqueu, e este era chefe dos publicanos, e era rico.
\VS{3}E procurava ver quem era Jesus, e não podia, por causa da multidão, pois era pequeno de altura.
\VS{4}E correndo com antecedência, subiu em uma árvore de frutos que parecem figos, para o ver; porque ele passaria por ali.
\VS{5}E quando Jesus chegou a aquele lugar, olhando para cima, o viu, e disse-lhe: Zaqueu, apressa-te, e desce; porque hoje é necessário que eu fique em tua casa.
\VS{6}E apressando-se, desceu, e o recebeu com alegria.
\VS{7}E todos, vendo
{\ADD{isto}} , murmuravam, dizendo: Ele entrou para se hospedar com um homem pecador.
\VS{8}E Zaqueu, levantando-se, disse ao Senhor: Senhor, eis que dou a metade de meus bens aos pobres; e se eu consegui algo enganando a alguém, eu
{\ADD{o}} devolvo quatro vezes mais.
\VS{9}E Jesus lhe disse: Hoje houve salvação nesta casa, porque ele também é filho de Abraão.
\VS{10}Porque o Filho do homem veio para buscar, e para salvar o que tinha se perdido.
\VS{11}E ouvindo eles estas coisas,
{\ADD{Jesus}} prosseguiu, e disse uma parábola, porque estava perto de Jerusalém, e pensavam que logo o Reino seria manifesto.
\VS{12}Disse, pois: Um certo homem nobre partiu para uma terra distante.
\VS{13}E chamando dez servos seus, deu-lhes dez minas, e disse-lhes: Investi até que eu venha;
\VS{14}E seus cidadãos o odiavam; e mandaram representantes depois dele, dizendo: Não queremos que este reine sobre nós.
\VS{15}E aconteceu que, quando ele voltou, tendo tomado o reino, disse que lhe chamassem a aqueles servos, a quem tinha dado o dinheiro, para saber o que cada um tinha ganho fazendo investimentos.
\VS{16}E veio o primeiro, dizendo: Senhor, tua mina rendeu outras dez minas.
\VS{17}E ele lhe disse: Ótimo, bom servo! Por teres sido fiel no pouco, terás autoridade sobre dez cidades.
\VS{18}E veio o segundo, dizendo: Senhor, tua mina rendeu cinco minas.
\VS{19}E disse também a este: E tu
{\ADD{governarás}} cinco cidades.
\VS{20}E veio outro, dizendo: Eis aqui tua mina, que guardei em um lenço.
\VS{21}Porque tive medo de ti, que és um homem rigoroso, que tomas o que não puseste, e colhes o que não semeaste.
\VS{22}Porém ele lhe disse: Servo mau, por tua boca eu te julgarei; tu sabias que eu era um homem rigoroso, que tomo o que não pus, e que colho o que não semeei;
\VS{23}Por que, então, não puseste meu dinheiro no banco; e quando eu viesse, o receberia de volta com juros?
\VS{24}E disse aos que estavam com ele: Tirai-lhe a mina, e dai-a ao que tem as dez minas.
\VS{25}E eles lhe disseram: Senhor, ele
{\ADD{já}} tem dez minas.
\VS{26}{\ADD{O senhor respondeu:}} Porque eu vos digo, que todo aquele que tiver, lhe será dado; mas ao que não tiver, até o que tem, lhe será tirado.
\VS{27}Porém a aqueles meus inimigos, que não quiseram que eu reinasse sobre eles, trazei
{\ADD{- os}} aqui, e matai
{\ADD{-os}} diante de mim.
\VS{28}E dito isto, ele foi caminhando adiante, subindo para Jerusalém.
\VS{29}E aconteceu que, chegando perto de Betfagé, e de Betânia, ao monte chamado das Oliveiras, mandou a dois de seus discípulos,
\VS{30}Dizendo: Ide à aldeia que está em frente; onde, ao entrardes, achareis um potro atado, em que ninguém jamais se sentou; soltai-o, e trazei
{\ADD{-o}} .
\VS{31}E se alguém vos perguntar: Por que
{\ADD{o}} soltais? Direis assim a ele: Porque o Senhor precisa dele.
\VS{32}E indo os que tinham sido mandados, acharam como lhes disse:
\VS{33}E soltando o potro, seus donos lhe disseram: Por que soltais o potro?
\VS{34}E eles disseram: O Senhor precisa dele.
\VS{35}E o trouxeram a Jesus; e lançando suas roupas sobre o potro, puseram Jesus montado nele.
\VS{36}E indo ele andando, estendiam suas roupas pelo caminho.
\VS{37}E quando já chegava perto da descida do monte das Oliveiras, toda a multidão de discípulos, com alegria começou a louvar a Deus em alta voz, por todas as maravilhas que tinham visto,
\VS{38}Dizendo: Bendito o Rei que vem em nome do Senhor; paz no céu, e glória nas alturas.
\VS{39}E alguns dos fariseus da multidão lhe disseram: Mestre, repreende a teus discípulos.
\VS{40}E respondendo ele, disse-lhes: Digo-vos, que se estes se calarem, as pedras clamariam.
\VS{41}E quando já estava chegando, viu a cidade, e chorou por causa dela,
\VS{42}Dizendo: Ah, se tu também conhecesses, pelo menos neste teu dia, aquilo que lhe traria paz! Mas agora
{\ADD{isto}} está escondido de teus olhos.
\VS{43}Porque dias virão sobre ti, em que teus inimigos lhe cercarão com barricadas, e ao redor te sitiarão, e lhe pressionarão por todos os lados.
\VS{44}E derrubarão a ti, e a teus filhos; e não deixarão em ti pedra sobre pedra, porque não conheceste o tempo em que foste visitada.
\VS{45}E entrando no Templo, começou a expulsar a todos os que vendiam e compravam ali,
\VS{46}Dizendo-lhes: Está escrito: Minha casa é casa de oração; Mas vós a tendes feito um esconderijo de ladrões.
\VS{47}E ensinava diariamente no Templo; e os chefes dos sacerdotes, e os escribas, e os chefes do povo, procuravam matá-lo.
\VS{48}E não achavam como fazer, porque todo o povo o ouvia com muita atenção.

\par }\Chap{20}{\PP \VerseOne{1}E aconteceu, num daqueles dias que, enquanto ele estava ensinando ao povo no Templo, e anunciando o Evangelho, vieram até ele os chefes dos sacerdotes, e os escribas com os anciãos.
\VS{2}E falaram-lhe, dizendo: Dize-nos, com que autoridade fazes estas coisas? Ou quem é o que te deu esta autoridade?
\VS{3}E respondendo ele, disse-lhes: Também eu vos perguntarei algo, e dizei-me:
\VS{4}O batismo de João era do céu, ou dos homens?
\VS{5}E eles discutiam entre si, dizendo: Se dissermos: Do céu, ele nos dirá: Por que, então, vós não
{\ADD{o}} crestes?
\VS{6}E se dissermos: Dos homens, todo o povo nos apedrejará; pois estão convencidos de que João era profeta.
\VS{7}E responderam que não sabiam de onde
{\ADD{era}} .
\VS{8}E Jesus lhes disse: Nem eu vos direi com que autoridade eu faço estas coisas.
\VS{9}E começou a dizer ao povo esta parábola: Um certo homem plantou uma vinha, e a arrendou a
{\ADD{uns}} lavradores, e viajou para outro país por muito tempo.
\VS{10}E certo tempo
{\ADD{depois}} mandou um servo aos lavradores, para que lhe dessem do fruto da vinha; mas os lavradores, espancando-o, mandaram
{\ADD{-no}} sem coisa alguma.
\VS{11}E voltou a mandar outro servo; mas eles, espancando e humilhando também
{\ADD{a ele}} ,o mandaram sem nada.
\VS{12}E voltou a mandar ao terceiro; mas eles, ferindo também a este,
{\ADD{o}} expulsaram.
\VS{13}E o senhor da vinha disse: Que farei? Mandarei a meu filho amado; talvez quando o verem,
{\ADD{o}} respeitarão.
\VS{14}Mas os lavradores, vendo-o, discutiram entre si, dizendo: Este é o herdeiro; vamos matá-lo, para que a herança venha a ser nossa.
\VS{15}E expulsando-o da vinha,
{\ADD{o}} mataram. O que, então, lhes fará o senhor da vinha?
\VS{16}Virá, e destruirá a estes lavradores, e dará a vinha a outros.E eles, ouvindo isto, disseram: Que
{\ADD{isto}} nunca aconteça!
\VS{17}Mas
{\ADD{Jesus}} ,olhando para eles, disse: Por que, então, isto está escrito: A pedra que os construtores rejeitaram, essa foi posta como a principal da esquina?
\VS{18}Todo aquele que cair sobre aquela pedra, se quebrará em pedaços; e aquele sobre quem ela cair, se fará pó.
\VS{19}E os chefes dos sacerdotes e os escribas queriam detê-lo naquela mesma hora, mas temiam ao povo; porque entenderam que foi contra eles que ele tinha dito a parábola.
\VS{20}E, observando-o, mandaram espiões, que fingissem ser justos, para o pegarem por meio de algo que ele dissesse, e o entregarem ao poder e autoridade do governador.
\VS{21}E perguntaram-lhe, dizendo: Mestre, nós sabemos que falas e ensinas corretamente, e que não te importas com as aparências, mas na verdade tu ensinas o caminho de Deus.
\VS{22}É lícito para nós dar tributo a César, ou não?
\VS{23}E ele, entendendo a astúcia deles, disse-lhes: Por que me tentais?
\VS{24}Mostrai-me uma moeda; ela tem a imagem e a inscrição de quem?E eles, respondendo, disseram: De César.
\VS{25}Então lhes disse: Dai pois a César o que
{\ADD{é}} de César, e a Deus o que
{\ADD{é}} de Deus.
\VS{26}E não puderam lhe pegar em algo que ele tenha dito diante do povo; e maravilhados de sua resposta, calaram-se.
\VS{27}E chegando-se alguns dos saduceus, que negam haver a ressurreição, perguntaram-lhe,
\VS{28}Dizendo: Mestre, Moisés nos escreveu, que se o irmão de alguém morrer, tendo mulher, e morrer sem filhos, o irmão deve tomar a mulher, e gerar descendência a seu irmão.
\VS{29}Houve, pois, sete irmãos, e o primeiro tomou mulher, e morreu sem filhos.
\VS{30}E o segundo a tomou; e
{\ADD{também}} este morreu sem filhos.
\VS{31}E o terceiro a tomou, e assim também os sete, e não deixaram filhos, e morreram.
\VS{32}E por fim, depois de todos, morreu também a mulher.
\VS{33}Na ressurreição, pois, ela será mulher de qual deles? Pois os sete a tiveram por mulher.
\VS{34}E respondendo Jesus, disse-lhes: Os filhos destes tempos se casam, e se dão em casamento.
\VS{35}Mas os que forem considerados dignos de alcançarem aqueles tempos futuros, e da ressurreição dos mortos, nem se casarão, nem se darão em casamento.
\VS{36}Porque já não podem mais morrer; pois são iguais aos anjos; e são filhos de Deus, pois são filhos da ressurreição.
\VS{37}E até Moisés mostrou, junto à sarça, que os mortos ressuscitam, quando ele chama ao Senhor de Deus de Abraão, e Deus de Isaque, e Deus de Jacó.
\VS{38}Ora, ele não é Deus de mortos, mas de vivos; pois todos vivem por causa dele.
\VS{39}E alguns dos escribas, respondendo, disseram: Mestre, bem disseste.
\VS{40}E não ousavam lhe perguntar mais nada.
\VS{41}E ele lhes disse: Como dizem que o Cristo é filho de Davi?
\VS{42}Pois o próprio Davi diz no livro dos Salmos: Disse o Senhor a meu Senhor: Senta-te à minha direita,
\VS{43}Até que eu ponha teus inimigos por escabelo de teus pés.
\VS{44}Se Davi o chama de Senhor, como, então, é seu filho?
\VS{45}Enquanto todo o povo estava ouvindo, ele disse a seus discípulos:
\VS{46}Tomai cuidado com os escribas, que querem andar roupas compridas, e amam as saudações nas praças, e as primeiras cadeiras nas sinagogas, e os primeiros assentos nos jantares.
\VS{47}Que devoram as casas das viúvas, e fingem fazer longas orações. Estes receberão mais grave condenação.

\par }\Chap{21}{\PP \VerseOne{1}E ele, olhando, viu os ricos lançarem suas ofertas na arca do tesouro
{\ADD{do templo}} .
\VS{2}E viu também uma pobre viúva lançar ali duas pequenas moedas.
\VS{3}E disse: Em verdade vos digo que esta pobre viúva lançou mais do que todos,
\VS{4}Porque todos aqueles lançaram para as ofertas de Deus daquilo que lhes sobrava; mas esta
{\ADD{viúva}} , de sua pobreza, lançou todo o sustento que tinha.
\VS{5}E alguns estavam falando do templo, que era adornado com formosas pedras e ofertas. Então
{\ADD{Jesus}} disse:
\VS{6}Destas coisas que vedes, dias virão, em que não se deixará pedra sobre pedra, que não seja derrubada.
\VS{7}E perguntaram-lhe, dizendo: Mestre, quando, pois, serão essas coisas? E que sinal haverá, quando essas coisas vierem a acontecer?
\VS{8}Então ele disse: Olhai para que não vos enganem, porque muitos virão em meu nome, dizendo: Eu sou
{\ADD{o Cristo}} . E o tempo já está perto; portanto não os sigais.
\VS{9}E quando ouvirdes de guerras, e de rebeliões, não vos espanteis. Porque é necessário que estas coisas aconteçam primeiro; mas ainda não é o fim.
\VS{10}Então lhe disse: Então se levantará nação contra nação, e reino contra reino.
\VS{11}E haverá em vários lugares grandes terremotos e fomes, e pragas; haverá também coisas espantosas, e grandes sinais do céu.
\VS{12}Mas antes de tudo isto, eles vos impedirão e vos perseguirão,
{\ADD{vos}} entregando em sinagogas e prisões,
{\ADD{e}} vos trazendo diante de reis, e governadores, por causa do meu nome.
\VS{13}E isto vos acontecerá para haver testemunho.
\VS{14}Portanto, que vós decidais nos vossos corações não planejar como direis em vossa defesa,
\VS{15}porque eu vos darei boca e sabedoria, para que todos os que forem contra vós não posam
{\ADD{vos}} contradizer ou resistir.
\VS{16}E vós sereis entregues até pelos pais, e irmãos, e parentes, e amigos; e
{\ADD{alguns}} de vós serão mortos.
\VS{17}E vós sereis odiados por todos por causa do meu nome.
\VS{18}Mas nem um cabelo de vossa cabeça parecerá.
\VS{19}Por vossa paciência ganhareis vossas almas.
\VS{20}Porém quando virdes a Jerusalém cercada de exércitos, sabei então, que próxima está sua desolação.
\VS{21}Então, os que estiverem na Judeia, fujam para os montes; e os que estiverem no meio dela, saiam, e os que estiverem nos campos, não entrem nela.
\VS{22}Porque estes são dias de vingança, para que todas as coisas que estão escritas se cumpram.
\VS{23}Mas ai das grávidas, e das que amamentarem naqueles dias; porque grande calamidade haverá na terra, e ira contra este povo.
\VS{24}E cairão pela lâmina de espada, e serão levados cativos para todas as nações; e Jerusalém será pisada pelos gentios, até que os tempos dos gentios se cumpram.
\VS{25}E haverá sinais no Sol,
{\ADD{na}} Lua, e
{\ADD{nas}} estrelas; e na terra sofrimento entre as nações, como o rugir e agitar do mar.
\VS{26}As pessoas desmaiarão de medo e da expectativa das coisas que vão acontecer ao mundo, porque os poderes dos céus serão abalados.
\VS{27}E então verão o Filho do homem vir em uma nuvem com grande poder e glória.
\VS{28}Ora, quando estas coisas começarem a acontecer, olhai para cima, e levantai vossas cabeças, porque vosso resgate está perto.
\VS{29}E disse-lhes uma parábola: Olhai a figueira, e todas as árvores;
\VS{30}Quando vós vedes elas já brotando, sabeis por vós mesmos que o verão já está perto.
\VS{31}Assim também vós, quando virdes acontecer estas coisas, sabei que o Reino de Deus já está perto.
\VS{32}Em verdade vos digo, que esta geração não passará, até que tudo aconteça.
\VS{33}O céu e a terra passarão, mas minhas palavras de maneira nenhuma passarão.
\VS{34}E olhai por vós, para que vossos corações não venham a se encher de ressaca e embriaguez, e das preocupações d
{\ADD{esta}} vida; e vos venha aquele dia de surpresa.
\VS{35}Porque virá como uma armadilha sobre todos os que habitam sobre a face de toda a terra.
\VS{36}Então vigiai sempre, orando para que sejais considerados dignos de escaparem de todas as coisas que irão acontecer, e de ficarem de pé diante do Filho do homem.
\VS{37}E ensinava durante os dias no Templo, porém, às noites saía e as passava no monte, chamado as Oliveiras.
\VS{38}E todo o povo vinha até ele de manhã cedo ao templo, para o ouvir.

\par }\Chap{22}{\PP \VerseOne{1}Estava perto a festa dos pães sem fermento, chamada de páscoa.
\VS{2}E os chefes dos sacerdotes, e os escribas procuravam um meio de o matar, pois eles temiam ao povo.
\VS{3}E Satanás entrou no Judas que era chamado Iscariotes, que era um dos doze.
\VS{4}E foi, e falou com os chefes dos sacerdotes e os oficiais, sobre como
{\ADD{o}} entregaria para eles.
\VS{5}E
{\ADD{estes}} se alegraram, e concordaram em lhe dar dinheiro.
\VS{6}E
{\ADD{lhes}} prometeu, e buscava oportunidade para o entregar quando não houvesse uma multidão.
\VS{7}E veio o dia dos
{\ADD{pães}} sem fermento, em que se devia fazer o sacrifício da páscoa.
\VS{8}E
{\ADD{Jesus}} mandou a Pedro, e a João, dizendo: Ide, preparai-nos a páscoa, para que
{\ADD{a}} comamos.
\VS{9}E eles lhe disseram: Onde queres que
{\ADD{a}} preparemos?
\VS{10}E ele lhes disse: Eis que, quando entrardes na cidade, um homem com um vaso de água vos encontrará; segui-o até a casa onde ele entrar.
\VS{11}E direis ao dono da casa: O Mestre te diz: Onde está o salão onde comerei a páscoa com meus discípulos?
\VS{12}Então ele vos mostrará um grande salão já arrumado; preparai-a ali.
\VS{13}E indo eles, acharam como lhes tinha dito; e prepararam a páscoa.
\VS{14}E vinda a hora, sentou-se
{\ADD{à mesa}} ,e com ele os doze apóstolos.
\VS{15}E disse-lhes: Muito desejei comer convosco esta páscoa, antes que eu sofra.
\VS{16}Porque eu vos digo, que dela não mais comerei, até que
{\ADD{isto}} se cumpra no Reino de Deus.
\VS{17}E tomando o copo, e tendo agradecido
{\ADD{a Deus}} ,disse: Tomai-o, e reparti
{\ADD{-o}} entre vós.
\VS{18}Porque vos digo, que do fruto da videira eu não beberei, até que o Reino de Deus venha.
\VS{19}E tomando o pão, e tendo agradecido
{\ADD{a Deus}} ,partiu-o, e o deu a eles, dizendo: Isto é o meu corpo, que é dado por vós; fazei isto em memória de mim.
\VS{20}De modo semelhante também com o copo, depois da ceia, disse: Este copo
{\ADD{é}} o Novo Testamento em meu sangue, que é derramado por vós.
\VS{21}Porém eis que a mão do que me trai
{\ADD{está}} comigo à mesa.
\VS{22}E realmente o Filho do homem vai conforme o que está determinado; mas ai daquele homem por quem é traído!
\VS{23}E começaram a perguntar entre si, qual deles seria o que faria isto.
\VS{24}E houve também uma briga entre eles, sobre qual deles era considerado o maior.
\VS{25}E
{\ADD{Jesus}} lhes disse: Os reis dos gentios os dominam, e os que exercem autoridade sobre eles são chamados de benfeitores;
\VS{26}Mas não
{\ADD{seja}} assim entre vós; antes o maior de vós seja como o menor; e o que lidera, como o que serve.
\VS{27}Porque qual é maior? O que se senta
{\ADD{à mesa}} ,ou o que serve? Por acaso não é o que se senta
{\ADD{à mesa}} ? Porém eu estou entre vós como aquele que serve.
\VS{28}E vós sois os que tendes permanecido comigo em minhas tentações.
\VS{29}E eu vos determino um Reino, assim como meu pai o determinou a mim.
\VS{30}Para que em meu Reino, comais e bebais à minha mesa; e vos senteis sobre tronos, julgando as doze tribos de Israel.
\VS{31}Disse também o Senhor: Simão, Simão; eis que Satanás vos pediu, para
{\ADD{vos}} peneirar como trigo;
\VS{32}Mas eu roguei por ti, que tua fé não se acabe; e quando tu te converteres, fortaleça teus irmãos.
\VS{33}E ele lhe disse: Senhor, estou preparado para ir contigo até à prisão, e à morte.
\VS{34}Mas ele disse: Pedro, eu te digo que hoje o galo não cantará, antes que me negues três vezes que me conheces.
\VS{35}E disse a eles: Quando vos mandei sem bolsa, e sem sacola, e sem sandálias, por acaso algo vos faltou?E disseram: Nada.
\VS{36}Então, ele lhes disse: Mas agora, quem tem bolsa, tome-a, como também a sacola; e o que não tem espada, venda sua roupa, e compre uma.
\VS{37}Porque eu vos digo, que ainda é necessário que se cumpra em mim aquilo que está escrito: E ele foi contado com os malfeitores. Porque aquilo que é sobre mim tem que se cumprir.
\VS{38}E eles disseram: Senhor, eis aqui duas espadas.E ele lhes disse: É o suficiente.
\VS{39}E saindo, foi, como de costume, para o monte das Oliveiras; e os seus discípulos também o seguiram.
\VS{40}E quando chegou a aquele lugar, disse-lhes: Orai para que não entreis em tentação.
\VS{41}E se afastou deles,
{\ADD{à distância}} de um tiro de pedra. E pondo-se de joelhos, orava,
\VS{42}Dizendo: Pai, se tu quiseres, passa este copo de mim; porém não se faça minha vontade, mas a tua.
\VS{43}E apareceu-lhe um anjo do céu, que o fortalecia.
\VS{44}E estando em angústia, orava mais intensamente. E seu suor se fez como gotas de sangue, que desciam até o chão.
\VS{45}E ele, levantando-se da oração, veio a seus discípulos, e os achou dormindo por causa da tristeza.
\VS{46}E disse-lhes: Por que estais dormindo? Levantai-vos, e orai, para que não entreis em tentação.
\VS{47}E enquanto ele ainda estava falando, eis que uma multidão
{\ADD{chegou}} ; e um dos doze, o que se chamava Judas, ia adiante deles, e se aproximou de Jesus, para o beijar.
\VS{48}E Jesus lhe disse: Judas, com um beijo trais ao Filho do homem?
\VS{49}E os que estavam com ele, vendo o que iria acontecer, disseram-lhe: Senhor, feriremos com a espada?
\VS{50}E um deles feriu a um servo do chefe dos sacerdotes, e cortou-lhe a orelha direita.
\VS{51}E respondendo Jesus, disse: Para com isto! E tocando-lhe a orelha, o curou.
\VS{52}E disse Jesus aos chefes dos sacerdotes, e aos oficiais do Templo, e aos anciãos, que tinham vindo contra ele: Como
{\ADD{se eu fosse}} ladrão, saístes com espadas e bastões?
\VS{53}Estando eu convosco todo dia no Templo, contra mim não me prendestes; mas esta é a vossa hora, e
{\ADD{sob}} a autoridade das trevas.
\VS{54}E prendendo-o,
{\ADD{o}} trouxeram e o puseram na casa do sumo sacerdote. E Pedro
{\ADD{o}} seguia de longe.
\VS{55}E acenderam fogo no meio do pátio, e sentaram-se juntos, e Pedro se sentou entre eles.
\VS{56}E uma serva, vendo-o sentado junto ao fogo, fixando o olhos nele, disse: Este também estava com ele.
\VS{57}Porém ele o negou, dizendo: Mulher, eu não o conheço.
\VS{58}E pouco depois, outro o viu, e disse: Também tu és um deles.
\VS{59}E quando já tinha passado quase uma hora, outro afirmava, dizendo: Verdadeiramente também este estava com ele, porque também é galileu.
\VS{60}E Pedro disse: Homem, não sei o que dizes.E logo, estando ele ainda falando, cantou o galo.
\VS{61}E o Senhor, virando-se, olhou para Pedro; e Pedro se lembrou da palavra do Senhor, como lhe tinha dito: Antes que o galo cante, tu me negarás três vezes.
\VS{62}E Pedro, saindo, chorou amargamente.
\VS{63}E os homens que tinham prendido a Jesus zombavam dele, ferindo-o;
\VS{64}E cobrindo-o, feriam-no no rosto; e perguntavam-lhe, dizendo: Profetiza, quem é o que te feriu?
\VS{65}E diziam muitas outras coisas contra ele, insultando-o.
\VS{66}E quando já era de dia, juntaram-se os anciãos do povo, e os chefes dos sacerdotes, e os escribas, e o trouxeram ao supremo conselho,
\FTNT{*}{{\FR 22:66 }
supremo conselho
lit. sinédrio – o mais importante conselho ou tribunal para os judeus
}
\VS{67}dizendo, Tu és o Cristo? Dize-nos. E ele lhes disse: Se eu vos disser, não o crereis.
\VS{68}E também se eu perguntar, não me respondereis, nem
{\ADD{me}} soltareis.
\VS{69}A partir de agora o Filho do homem se sentará à direita do poder de Deus.
\VS{70}E todos disseram: Então tu és o Filho de Deus?E ele lhes disse: Vós dizeis que eu sou.
\VS{71}E eles disseram: Para que precisamos de mais testemunho? Pois nós mesmos o ouvimos de sua boca.

\par }\Chap{23}{\PP \VerseOne{1}E levantando-se toda a multidão deles, o levaram a Pilatos.
\VS{2}E começaram a acusá-lo, dizendo: Encontramos este
{\ADD{homem}} ,que perverte a nação, e proíbe dar tributo a César, dizendo que ele mesmo é Cristo, o Rei.
\VS{3}E Pilatos lhe perguntou, dizendo: Tu és o Rei dos judeus?E respondendo, ele lhe disse: Tu o dizes.
\VS{4}E Pilatos disse aos chefes dos sacerdotes, e às multidões: Não acho culpa nenhuma neste homem.
\VS{5}Mas eles insistiam, dizendo: Ele incita ao povo, ensinando por toda a Judeia, começando desde a Galileia até aqui.
\VS{6}Então Pilatos, ouvindo
{\ADD{falar}} da Galileia, perguntou se aquele homem era galileu.
\VS{7}E quando soube que era da jurisdição de Herodes, ele o entregou a Herodes, que naqueles dias também estava em Jerusalém.
\VS{8}E Herodes, ao ver Jesus, alegrou-se muito, porque havia muito tempo que desejava o ver, pois ouvia muitas coisas sobre ele; e esperava ver algum sinal feito por ele.
\VS{9}E perguntava-lhe com muitas palavras, mas ele nada lhe respondia;
\VS{10}E estavam
{\ADD{lá}} os chefes dos sacerdotes, e os escribas, acusando-o com veemência.
\VS{11}E Herodes, com seus soldados, desprezando-o, e escarnecendo dele, o vestiu com uma roupa luxuosa, e o enviou de volta a Pilatos.
\VS{12}E no mesmo dia Pilatos e Herodes se fizeram amigos; porque antes tinham inimizade um contra o outro.
\VS{13}E Pilatos, convocando aos chefes dos sacerdotes, aos líderes, e ao povo, disse-lhes:
\VS{14}Vós me trouxestes a este homem, como que perverte o povo; e eis que eu, examinando-o em vossa presença, nenhuma culpa eu acho neste homem, das de que o acusais.
\VS{15}E nem também Herodes; porque a ele eu vos remeti; e eis que ele nada fez para que seja digno de morte.
\VS{16}Então eu o castigarei, e
{\ADD{depois}} o soltarei.
\VS{17}E ele tinha de soltar-lhes alguém durante a festa.
\VS{18}Porém todos clamavam juntos, dizendo: Tirai-o daqui! E soltai Barrabás para nós!
\VS{19}(O qual por uma rebelião feita na cidade, e
{\ADD{por}} uma morte, tinha sido lançado na prisão).
\VS{20}Pilatos falou-lhes então outra vez, querendo soltar Jesus.
\VS{21}Mas eles clamavam, dizendo: Crucifica
{\ADD{-o}} ! crucifica-o!
\VS{22}E ele lhes disse a terceira vez: Pois que mal este fez? Nenhuma culpa de morte nele eu achei. Então eu o castigarei, e
{\ADD{depois}} o soltarei.
\VS{23}Mas eles continuavam, com grandes gritos, pedindo que fosse crucificado. E seus gritos, e
{\ADD{os}} dos chefes dos sacerdotes, prevaleceram.
\VS{24}Então Pilatos julgou que se fizesse o que pediam.
\VS{25}E soltou-lhes ao que fora lançado na prisão por uma rebelião e
{\ADD{uma}} morte, que
{\ADD{era}} o que pediam; porém a Jesus
{\ADD{lhes}} entregou à sua vontade.
\VS{26}E enquanto o levavam, tomaram a um Simão Cireneu, que vinha do campo, e puseram-lhe a cruz às costas, para que a levasse atrás de Jesus.
\VS{27}E seguia-o uma grande multidão do povo, e de mulheres, as quais também ficavam desconsoladas, e lamentavam por ele.
\VS{28}E Jesus, virando-se para elas, disse: Filhas de Jerusalém, não choreis por mim, mas chorai por vós mesmas, e por vossos filhos.
\VS{29}Porque eis que vêm dias em que dirão: Bem-aventuradas as estéreis, e os ventres que não deram a luz, e os peitos que não amamentaram.
\VS{30}Então começarão a dizer aos montes: Cai sobre nós; E aos morros: Cobri-nos!
\VS{31}Porque, se fazem isto à árvore verde, o que se fará com a
{\ADD{árvore}} seca?
\VS{32}E também levaram outros dois, que eram malfeitores, para matar com ele.
\VS{33}E quando chegaram ao lugar, chamado a Caveira, o crucificaram ali, e aos malfeitores, um à direita, e outro à esquerda.
\VS{34}E Jesus dizia: Pai, perdoa-lhes, porque não sabem o que fazem. E, repartindo suas roupas, lançaram sortes.
\VS{35}E o povo estava olhando; e os líderes também zombavam com eles, dizendo: Salvou a outros, salve agora a si mesmo, se é o Cristo, o escolhido de Deus.
\VS{36}E os soldados também escarneciam dele, aproximando-se dele, e mostrando-lhe vinagre;
\VS{37}E dizendo: Se tu és o Rei dos judeus, salva a ti mesmo.
\VS{38}E também estava acima dele um título escrito com letras gregas, romanas e hebraicas: ESTE É O REI DOS JUDEUS.
\VS{39}E um dos malfeitores que estavam pendurados o insultava, dizendo: Se tu és o Cristo, salva a ti mesmo, e a nós.
\VS{40}Porém o outro, respondendo, repreendia-o, dizendo: Tu ainda não temes a Deus,
{\ADD{mesmo}} estando na mesma condenação?
\VS{41}E nós realmente
{\ADD{estamos sendo punidos}} justamente, porque estamos recebendo de volta merecidamente por aquilo que praticamos; mas este nada fez de errado.
\VS{42}E disse a Jesus: Senhor, lembra-te de mim, quando chegares em teu Reino.
\VS{43}E Jesus lhe disse: Em verdade te digo, hoje estarás comigo no paraíso.
\VS{44}E era já quase a hora sexta, e houve trevas em toda a terra, até a hora nona.
\VS{45}E o sol se escureceu, e o véu do templo se rasgou ao meio.
\VS{46}E Jesus, clamando em alta voz, disse: Pai, em tuas mãos eu entrego meu espírito.E tendo dito isto, parou de respirar.
\VS{47}E o centurião, vendo o o que tinha acontecido, deu glória a Deus, dizendo: Verdadeiramente este homem era justo.
\VS{48}E todas as multidões que se juntavam para observar, vendo o que tinha acontecido, voltaram, batendo nos peitos.
\VS{49}E todos os seus conhecidos, e as mulheres que acompanhando
{\ADD{-o}} desde a Galileia, tinham o seguido, estavam longe, vendo estas coisas.
\VS{50}E eis que um homem, de nome José, membro do conselho
{\ADD{de justiça}} ,sendo homem bom e justo.
\VS{51}(Que não tinha concordado, nem com o conselho, nem em atos que fizeram), da cidade de Arimateia,
{\ADD{da terra}} dos judeus, e que também esperava pelo Reino de Deus.
\VS{52}Este, chegando a Pilatos, pediu o corpo de Jesus.
\VS{53}E tendo o tirado, o envolveu em um tecido de linho, e o pôs em um sepulcro, escavado em uma rocha, onde nunca ainda tinha sido posto.
\VS{54}E era o dia da preparação, e o sábado estava começando.
\VS{55}E as mulheres que vieram com ele da Galileia também
{\ADD{o}} seguiram, e viram o sepulcro, e como seu corpo foi posto.
\VS{56}E elas, ao voltarem, prepararam materiais aromáticos e óleos perfumados. E descansaram o sábado, conforme o mandamento.

\par }\Chap{24}{\PP \VerseOne{1}E no primeiro
{\ADD{dia}} da semana, de madrugada bem cedo, foram ao sepulcro, levando
{\ADD{consigo}} os materiais aromáticos que tinham preparado; e algumas
{\ADD{outras}} junto delas.
\VS{2}E acharam a pedra já revolvida do sepulcro.
\VS{3}E entrando, não acharam o corpo do Senhor Jesus.
\VS{4}E aconteceu, que estando elas perplexas, eis que dois homens apareceram junto a elas, com roupas luminosas.
\VS{5}E estando elas com muito medo, e abaixando o rosto para o chão, eles lhes disseram: Por que buscam entre os mortos aquele que vive?
\VS{6}Ele não está aqui, mas já ressuscitou. Lembrai-vos de como ele vos falou, quando ainda estava na Galileia,
\VS{7}Dizendo: É necessário que o Filho do homem seja entregue nas mãos de homens pecadores, e
{\ADD{que}} seja crucificado, e ressuscite ao terceiro dia.
\VS{8}E se lembraram das palavras dele.
\VS{9}E, voltando do sepulcro, anunciaram todas estas coisas aos onze, e a todos os outros.
\VS{10}E eram Maria Madalena, e Joana, e Maria
{\ADD{mãe}} de Tiago, e as outras
{\ADD{que estavam}} com elas, que diziam estas coisas aos apóstolos.
\VS{11}E para eles, as palavras delas pareciam não ter sentido; e não creram nelas.
\VS{12}Porém Pedro, levantando-se, correu ao sepulcro; e abaixando-se, viu os tecidos postos separadamente; e saiu maravilhado com o que tinha acontecido.
\VS{13}E eis que dois deles iam naquele mesmo dia a uma aldeia, cujo nome era Emaús, que estava a sessenta estádios
{\ADD{de distância}} de Jerusalém.
\VS{14}E iam falando entre si de todas aquelas coisas que tinham acontecido.
\VS{15}E aconteceu que, enquanto eles estavam conversando entre si, e perguntando um ao outro, Jesus se aproximou, e foi junto deles.
\VS{16}Mas seus olhos foram retidos, para que não o reconhecessem.
\VS{17}E disse-lhes: Que conversas são essas, que vós discutis enquanto andam, e ficais tristes?
\VS{18}E um
{\ADD{deles}} ,cujo nome era Cleofas, respondendo-o, disse-lhe: És tu o único viajante em Jerusalém que não sabe as coisas que nela tem acontecido nestes dias?
\VS{19}E ele lhes disse: Quais?E eles lhe disseram: As sobre Jesus de Nazaré, a qual foi um homem profeta, poderoso em obras e em palavras, diante de Deus, e de todo o povo.
\VS{20}E como os chefes dos sacerdotes, e nossos líderes o entregaram à condenação de morte, e o crucificaram.
\VS{21}E nós esperávamos que ele fosse aquele que libertar a Israel; porém além de tudo isto, hoje é o terceiro dia desde que estas coisas aconteceram.
\VS{22}Ainda que também algumas mulheres dentre nós nos deixaram surpresos, as quais de madrugada foram ao sepulcro;
\VS{23}E não achando seu corpo, vieram, dizendo que também tinham visto uma aparição de anjos, que disseram que ele vive.
\VS{24}E alguns do que estão conosco foram ao sepulcro, e
{\ADD{o}} acharam assim como as mulheres tinham dito; porém não o viram.
\VS{25}E ele lhes disse: Ó tolos, e que demoram no coração para crerem em tudo o que os profetas falaram!
\VS{26}Por acaso não era necessário que o Cristo sofresse estas coisas, e
{\ADD{então}} entrar em sua glória?
\VS{27}E começando de Moisés, e por todos os profetas, lhes declarava em todas as Escrituras o que estava
{\ADD{escrito}} sobre ele.
\VS{28}E chegaram à aldeia para onde estavam indo; e ele agiu como se fosse para
{\ADD{um lugar}} mais distante.
\VS{29}E eles lhe rogaram, dizendo: Fica conosco, porque já é tarde, e o dia está entardecendo;E ele entrou para ficar com eles.
\VS{30}E aconteceu que, estando sentado com eles
{\ADD{à mesa}} ,tomou o pão, o benzeu,
{\ADD{o}} partiu, e o deu a eles.
\VS{31}E os olhos deles se abriram, e o reconheceram, e ele lhes desapareceu.
\VS{32}E diziam um ao outro: Por acaso não estava nosso coração ardendo em nós, quando ele falava conosco pelo caminho, e quando nos desvendava as Escrituras?
\VS{33}E levantando-se na mesma hora, voltaram para Jerusalém, e acharam reunidos aos onze, e aos que estavam com eles,
\VS{34}Que diziam: Verdadeiramente o Senhor ressuscitou, e já apareceu a Simão.
\VS{35}E eles contaram as coisas que
{\ADD{lhes aconteceram}} no caminho; e como foi reconhecido por eles quando partiu o pão.
\VS{36}E enquanto eles falavam disto, o próprio Jesus se pôs no meio deles, e lhes disse: Paz
{\ADD{seja}} convosco.
\VS{37}E eles, espantados, e muito atemorizados, pensavam que viam algum espírito.
\VS{38}E ele lhes disse: Por que estais perturbados, e por que sobem dúvidas em vossos corações?
\VS{39}Vede minhas mãos, e os meus pés, que sou em mesmo. Tocai-me, e vede, porque um espírito não tem carne nem ossos, como vós vedes que eu tenho.
\VS{40}E dizendo isto, lhes mostrou as mãos e os pés.
\VS{41}E eles, não crendo ainda, por causa da alegria, e maravilhados,
{\ADD{Jesus}} disse-lhes: Tendes aqui alguma coisa para comer?
\VS{42}Então eles lhe apresentaram parte de um peixe assado e de um favo de mel.
\VS{43}Ele pegou, e comeu diante deles.
\VS{44}E disse-lhes: Estas são as palavras que eu vos disse, enquanto ainda estava convosco, que era necessário que se cumprissem todas as coisas que estão escritas sobre mim na Lei de Moisés,
{\ADD{nos}} profetas, e
{\ADD{nos}} Salmos.
\VS{45}Então ele lhes abriu o entendimento, para que entendessem as Escrituras.
\VS{46}E disse-lhes: Assim está escrito, e assim era necessário que o Cristo sofresse, e que ao terceiro dia ressuscitasse dos mortos;
\VS{47}E que em seu nome fosse pregado arrependimento e perdão de pecados em todas as nações, começando de Jerusalém.
\VS{48}E destas coisas vós sois testemunhas.
\VS{49}E eis que eu envio a promessa de meu Pai sobre vós; porém ficai vós na cidade de Jerusalém, até que vos seja dado poder do alto.
\VS{50}E os levou para fora até Betânia, e levantando suas mãos, os abençoou.
\VS{51}E aconteceu que, enquanto os abençoava, ele se afastou deles, e foi conduzido para cima ao céu.
\VS{52}E eles, adorando-o, voltaram para Jerusalém com grande alegria;
\VS{53}E estavam sempre no Templo, louvando e bendizendo a Deus. Amém.\par }